%%%%%%%%%%%%%%%%%%%%%%%%%%%%%%%%%%%%%%%%%%%%%%%%
%%%%%%%%%%%%%%%%%%%%%%%%%%%%%%%%%%%%%%%%%%%%%%%%
%%%  Apart from this header, this file is    %%%
%%%  identical to PoHTC8.tex.  This file was %%%
%%%  posted to the arXiv on 7/9-2026.        %%%
%%%%%%%%%%%%%%%%%%%%%%%%%%%%%%%%%%%%%%%%%%%%%%%%
%%%%%%%%%%%%%%%%%%%%%%%%%%%%%%%%%%%%%%%%%%%%%%%%
\documentclass[12pt,reqno]{amsart}

\usepackage{geometry}
%\geometry{a4paper,margin=4.27382cm}
\geometry{a4paper,margin=2.0cm}
%\geometry{a4paper,left=2cm,right=2cm,top=2.9cm,bottom=2.3cm}
%\geometry{a4paper,left=1.2cm,right=1.2cm,top=1.55cm,bottom=1.55cm}


\raggedbottom


%Package hhline makes available \hhline for the tabular environment.
\usepackage{hhline}

%Package mathtools makes available \begin{psmallmatrix} ... \end{psmallmatrix}
\usepackage{mathtools}

\usepackage{mathdots}

%Enables coloured text.
\usepackage{color}

%Package pdfsync enables SHIFT+COMMAND+click syncing between the pdf
%file in Skim and the tex file in Aquamacs
%
\usepackage{pdfsync}

%Package enumitem makes possible the syntax
%\begin{itemize}[leftmargin=150pt]
%\item
% item
%\end{itemize}
\usepackage{enumitem}

%Package wasysym contains things like \pentagon, \hexagon, \varhexagon, \octagon.
\usepackage{wasysym}


%Package amssymb contains many nice mathematical symbols
%
\usepackage{amssymb}


%\usepackage{MnSymbol}
%Some say MnSymbol is incompatible with amssymb.

%\usepackage{fdsymbol}
%fdsymbol seems to be compatible with amssymb and gives access to many symbols, e.g. \cupdot for disjoint unions.  But it redefines many symbols to use an ugly font, so I prefer not to use it.

%Uncomment the following line to get nice mathcal letters, by way of
%the command \mathcal{some upper case letters}
%
%\usepackage{eucal}


%Uncomment the following line to get even nicer mathcal letters,
%by way of the command \mathscr{some upper case letters}
%
\usepackage{mathrsfs}


%The following package gives access to lower case blackboard bold,
%by way of the command \mathbbm{some letters}.
%
\usepackage{bbm}


%Uncomment the following line to get different(?) mathfrak letters
%
%\usepackage{eufrak}


%Uncomment the following line to get an alternative sort of script
%letters, by way of the command \mathpzc{some letters}
%
\DeclareMathAlphabet{\mathpzc}{OT1}{pzc}{m}{it}


%Package xy gives access to XYpic, a diagram environment
%
\usepackage[all]{xy}

%Package tikz gives access to TikZ, another diagram environment
%There seem to be many options for loading libraries; one needs to 
%consult the TikZ manual
%
\usepackage{tikz}
\usetikzlibrary{arrows,matrix,decorations.pathmorphing,decorations.pathreplacing,positioning,shapes.geometric,shapes.misc,decorations.markings,decorations.fractals,calc,patterns}


%Package tikz-cd is for commutative diagrams.
\usepackage{tikz-cd}


\usepackage{graphicx}
%Package graphicx permits commands like \rotatebox{90}{\subseteq}


%\usepackage{placeins}
%This package makes it possible to use the command \FloatBarrier
%in the document.  It stops figures moving across the barrier.


\usepackage{float}
%This package adds the possibility to write
%\begin{figure}[H]
%...
%\end{figure}
%and get the figure placed precisely there in the text ("H" for "Here").


%Package hyperref turns references into links.
%
%According to online folklore, the packages
%float, hyperref, algorithm
%must be loaded in that order (at the moment
%I don't use algorithm).
%
%\usepackage{hyperref}


\usepackage[bottom]{footmisc}
%This avoids figures being placed below footnotes.


\usepackage{moreenum}
%This package enables Greek letters as enumerate labels by the syntax:
%\begin{enumerate}[label=(\greek*)]
%...
%\end{enumerate}


\usepackage{makecell}
%This package makes it possible to split a cell diagonally in the tabular environment 
%using the command \diaghead.


%\usepackage{quiver}
%This is included by advice from the online diagram editor "quiver".


%The following line gives better baseline alignment in XYpic diagrams
%
\entrymodifiers={+!!<0pt,\fontdimen22\textfont2>}


%The following lines make the commands "dt" and "ddt" which can be used in maths mode thus to place fat dots and double dots above a letter: \dt{x} \ddt{x}
%\usepackage{accents}
%\newcommand*{\dt}[1]{%
%   \accentset{\mbox{\large\bfseries .}}{#1}}
%\newcommand*{\ddt}[1]{%
%   \accentset{\mbox{\large\bfseries .\hspace{-0.25ex}.}}{#1}}

%The following lines are the alternative where "dt" and "ddt" make ordinary dots.
\newcommand{\dt}{\dot}
\newcommand{\ddt}{\ddot}


%The following lines make the commands "newhat" and "newcheck" which can be used in maths mode thus: \newhat{c} \newcheck{c}
\usepackage{scalerel,stackengine}
\stackMath
\newcommand\newcheck[1]{%
\savestack{\tmpbox}{\stretchto{%
  \scaleto{%
    \scalerel*[\widthof{\ensuremath{#1}}]{\kern-.6pt\bigwedge\kern-.6pt}%
    {\rule[-\textheight/2]{1ex}{\textheight}}%WIDTH-LIMITED BIG WEDGE
  }{\textheight}% 
}{0.5ex}}%
\stackon[1pt]{#1}{\scalebox{-1}{\tmpbox}}%
}
\stackMath
\newcommand\newhat[1]{%
\savestack{\tmpbox}{\stretchto{%
  \scaleto{%
    \scalerel*[\widthof{\ensuremath{#1}}]{\kern-.6pt\bigwedge\kern-.6pt}%
    {\rule[-\textheight/2]{1ex}{\textheight}}%WIDTH-LIMITED BIG WEDGE
  }{\textheight}% 
}{0.5ex}}%
\stackon[1pt]{#1}{\scalebox{1}{\tmpbox}}%
}


%The following lines makes it possible to doctor the margins.  They are now obsolete because package geometry is used instead.
%\setlength{\textwidth}{172mm}
%The default \textheight is (close to) 206mm.
%\setlength{\textheight}{235mm}
%\addtolength{\oddsidemargin}{-2.2cm}
%\addtolength{\evensidemargin}{-2.2cm}
%\addtolength{\topmargin}{-11mm}


%Font for small sans serif letters
\font\sss=cmss8


%The following defines shorthands for upper case calligraphic letters
%
\def\cA{\mathscr{A}}
\def\cB{\mathscr{B}}
\def\cC{\mathscr{C}}
\def\cD{\mathscr{D}}
\def\cE{\mathscr{E}}
\def\cF{\mathscr{F}}
\def\cG{\mathscr{G}}
\def\cH{\mathscr{H}}
\def\cI{\mathscr{I}}
\def\cJ{\mathscr{J}}
\def\cK{\mathscr{K}}
\def\cL{\mathscr{L}}
\def\cM{\mathscr{M}}
\def\cN{\mathscr{N}}
\def\cO{\mathscr{O}}
\def\cP{\mathscr{P}}
\def\cQ{\mathscr{Q}}
\def\cR{\mathscr{R}}
\def\cS{\mathscr{S}}
\def\cT{\mathscr{T}}
\def\cU{\mathscr{U}}
\def\cV{\mathscr{V}}
\def\cW{\mathscr{W}}
\def\cX{\mathscr{X}}
\def\cY{\mathscr{Y}}
\def\cZ{\mathscr{Z}}


%The following defines shorthands for upper case blackboard bold letters
%
\def\BA{\mathbb{A}}
\def\BB{\mathbb{B}}
\def\BC{\mathbb{C}}
\def\BD{\mathbb{D}}
\def\BE{\mathbb{E}}
\def\BF{\mathbb{F}}
\def\BG{\mathbb{G}}
\def\BH{\mathbb{H}}
\def\BI{\mathbb{I}}
\def\BJ{\mathbb{J}}
\def\BK{\mathbb{K}}
\def\BL{\mathbb{L}}
\def\BM{\mathbb{M}}
\def\BN{\mathbb{N}}
\def\BO{\mathbb{O}}
\def\BP{\mathbb{P}}
\def\BQ{\mathbb{Q}}
\def\BR{\mathbb{R}}
\def\BS{\mathbb{S}}
\def\BT{\mathbb{T}}
\def\BU{\mathbb{U}}
\def\BV{\mathbb{V}}
\def\BW{\mathbb{W}}
\def\BX{\mathbb{X}}
\def\BY{\mathbb{Y}}
\def\BZ{\mathbb{Z}}


%The following defines shorthands for lower case blackboard bold letters
%The line  \usepackage{bbm}  above should be uncommented first
%
% \def\Ba{\mathbbm{a}}
% \def\Bb{\mathbbm{b}}
% \def\Bc{\mathbbm{c}}
% \def\Bd{\mathbbm{d}}
% \def\Be{\mathbbm{e}}
% \def\Bf{\mathbbm{f}}
% \def\Bg{\mathbbm{g}}
% \def\Bh{\mathbbm{h}}
% \def\Bi{\mathbbm{i}}
% \def\Bj{\mathbbm{j}}
\def\Bk{\mathbbm{k}}
% \def\Bl{\mathbbm{l}}
% \def\Bm{\mathbbm{m}}
% \def\Bn{\mathbbm{n}}
% \def\Bo{\mathbbm{o}}
% \def\Bp{\mathbbm{p}}
% \def\Bq{\mathbbm{q}}
% \def\Br{\mathbbm{r}}
% \def\Bs{\mathbbm{s}}
% \def\Bt{\mathbbm{t}}
% \def\Bu{\mathbbm{u}}
% \def\Bv{\mathbbm{v}}
% \def\Bw{\mathbbm{w}}
% \def\Bx{\mathbbm{x}}
% \def\By{\mathbbm{y}}
% \def\Bz{\mathbbm{z}}


%The following defines shorthands for upper case fraktur letters
%
\def\fA{\mathfrak{A}}
\def\fB{\mathfrak{B}}
\def\fC{\mathfrak{C}}
\def\fD{\mathfrak{D}}
\def\fE{\mathfrak{E}}
\def\fF{\mathfrak{F}}
\def\fG{\mathfrak{G}}
\def\fH{\mathfrak{H}}
\def\fI{\mathfrak{I}}
\def\fJ{\mathfrak{J}}
\def\fK{\mathfrak{K}}
\def\fL{\mathfrak{L}}
\def\fM{\mathfrak{M}}
\def\fN{\mathfrak{N}}
\def\fO{\mathfrak{O}}
\def\fP{\mathfrak{P}}
\def\fQ{\mathfrak{Q}}
\def\fR{\mathfrak{R}}
\def\fS{\mathfrak{S}}
\def\fT{\mathfrak{T}}
\def\fU{\mathfrak{U}}
\def\fV{\mathfrak{V}}
\def\fW{\mathfrak{W}}
\def\fX{\mathfrak{X}}
\def\fY{\mathfrak{Y}}
\def\fZ{\mathfrak{Z}}


%The following defines shorthands for lower case fraktur letters
%
\def\fa{\mathfrak{a}}
\def\fb{\mathfrak{b}}
\def\fc{\mathfrak{c}}
\def\frakd{\mathfrak{d}}
\def\fe{\mathfrak{e}}
\def\ff{\mathfrak{f}}
\def\frakg{\mathfrak{g}}
\def\fh{\mathfrak{h}}
\def\fraki{\mathfrak{i}}
\def\fj{\mathfrak{j}}
\def\fk{\mathfrak{k}}
%\def\fl{\mathfrak{l}}
\def\fm{\mathfrak{m}}
\def\fn{\mathfrak{n}}
\def\fo{\mathfrak{o}}
\def\fp{\mathfrak{p}}
\def\fq{\mathfrak{q}}
\def\fr{\mathfrak{r}}
\def\fs{\mathfrak{s}}
\def\ft{\mathfrak{t}}
\def\fu{\mathfrak{u}}
\def\fv{\mathfrak{v}}
\def\fw{\mathfrak{w}}
\def\fx{\mathfrak{x}}
\def\fy{\mathfrak{y}}
\def\fz{\mathfrak{z}}


%The following defines shorthands for upper case sans serif letters
%
\def\sA{\mathsf{A}}
\def\sB{\mathsf{B}}
\def\sC{\mathsf{C}}
\def\sD{\mathsf{D}}
\def\sE{\mathsf{E}}
\def\sF{\mathsf{F}}
\def\sG{\mathsf{G}}
\def\sH{\mathsf{H}}
\def\sI{\mathsf{I}}
\def\sJ{\mathsf{J}}
\def\sK{\mathsf{K}}
\def\sL{\mathsf{L}}
\def\sM{\mathsf{M}}
\def\sN{\mathsf{N}}
\def\sO{\mathsf{O}}
\def\sP{\mathsf{P}}
\def\sQ{\mathsf{Q}}
\def\sR{\mathsf{R}}
\def\sS{\mathsf{S}}
\def\sT{\mathsf{T}}
\def\sU{\mathsf{U}}
\def\sV{\mathsf{V}}
\def\sW{\mathsf{W}}
\def\sX{\mathsf{X}}
\def\sY{\mathsf{Y}}
\def\sZ{\mathsf{Z}}


%The following defines shorthands for lower case sans serif letters
%
\def\sa{\mathsf{a}}
\def\sb{\mathsf{b}}
\def\serifc{\mathsf{c}}
\def\sd{\mathsf{d}}
\def\se{\mathsf{e}}
\def\seriff{\mathsf{f}}
\def\sg{\mathsf{g}}
\def\sh{\mathsf{h}}
\def\si{\mathsf{i}}
\def\sj{\mathsf{j}}
\def\sk{\mathsf{k}}
\def\serifl{\mathsf{l}}
\def\sm{\mathsf{m}}
\def\sn{\mathsf{n}}
\def\so{\mathsf{o}}
\def\sp{\mathsf{p}}
\def\sq{\mathsf{q}}
\def\sr{\mathsf{r}}
%\def\ss{\mathsf{s}}
\def\st{\mathsf{t}}
\def\su{\mathsf{u}}
\def\sv{\mathsf{v}}
\def\sw{\mathsf{w}}
\def\sx{\mathsf{x}}
\def\sy{\mathsf{y}}
\def\sz{\mathsf{z}}


%\mhyphen will typeset a hyphen in maths mode.  Supposedly, it changes size and font depending on the current settings in maths mode.
\mathchardef\mhyphen="2D


%The following is Peter's mathematical operators
%
\def\ac{\operatorname{ac}}
\def\add{\operatorname{add}}
\def\Add{\operatorname{Add}}
\def\addgen{\operatorname{addgen}}
%The following line makes ascending dots for use in AR quivers
\def\adots{\mathinner{\mkern1mu\raise1.0pt\vbox{\kern7.0pt\hbox{.}}\mkern2mu\raise5.0pt\hbox{.}\mkern2mu\raise9.0pt\hbox{.}\mkern1mu}}
\def\amp{\operatorname{amp}}
\def\Ann{\operatorname{Ann}}
\def\ann{\operatorname{ann}}
\def\ast{{\textstyle *}}
\def\astsmall{{\scriptstyle *}}
\def\Aut{\operatorname{Aut}}
\def\b{\operatorname{b}}
\def\B{\operatorname{B}}
%\def\c{\operatorname{c}}
%\def\C{\operatorname{C}}
\def\cd{\operatorname{cd}}
\def\Cell{\operatorname{Cell}}
\def\Ch{\operatorname{Ch}}
\def\cls{\operatorname{cls}}
\def\CMreg{\operatorname{CMreg}}
\def\cocone{\operatorname{cocone}}
\def\coind{\operatorname{coind}}
\def\Coker{\operatorname{Coker}}
\def\colim{\operatorname{colim}}
\def\comp{\operatorname{comp}}
\newcommand\ConditionELong[4]{\mbox{$#3( #4,\Sigma^{ #1...#2 }#4 ) = 0$}}
%{\mbox{$\Ext^{#1...#2}_{#3}(#4,#4)=0$}}
\newcommand\ConditionEShort[3]
{\mbox{$#2( #3,\Sigma^{ #1 }#3 ) = 0$}}
%${\mbox{$\Ext^{#1}_{#2}(#3,#3)=0$}}
\def\cone{\operatorname{cone}}
\def\corad{\operatorname{corad}}
\def\Coslice{\operatorname{Coslice}}
\def\Cot{\operatorname{Cot}}
\newcommand{\cupdot}{\mathbin{\mathaccent\cdot\cup}}
\def\D{\cD}
\def\Dc{\D^{\operatorname{c}}}
\def\diag{\operatorname{diag}}
\def\dddots{\mathinner{\mkern1mu\raise10.0pt\vbox{\kern7.0pt\hbox{.}}\mkern2mu\raise5.3pt\hbox{.}\mkern2mu\raise1.0pt\hbox{.}\mkern1mu}}
\def\dddotssmall{\mathinner{\mkern1mu\raise7.0pt\vbox{\kern7.0pt\hbox{.}}\mkern-1mu\raise4pt\hbox{.}\mkern-1mu\raise1.0pt\hbox{.}\mkern1mu}}
\def\deg{\operatorname{deg}}
\def\depth{\operatorname{depth}}
\def\Db{\cD^{\operatorname{b}}}
\def\Df{\cD^{\operatorname{f}}}
\def\Diff{\operatorname{Diff}}
\def\dim{\operatorname{dim}}
\def\dirlim{\operatorname{lim}_{_{_{_{\hspace{-17pt} 
                                      \textstyle \longrightarrow}}}}}
\def\Dm{\D^-}
\def\Dp{\D^+}
\def\dtors{d\mbox{-tors}}
\def\dual{\operatorname{D}}
\def\End{\operatorname{End}}
\def\EssIm{\operatorname{Ess.\!Im}}
\def\Ext{\operatorname{Ext}}
\def\Extreg{\operatorname{Ext.\!reg}}
\def\Fac{\operatorname{Fac}}
\def\fd{\operatorname{fd}}
\def\fg{\operatorname{fg}}
\def\fin{\mathsf{fin}}
\def\finpres{\mathsf{fp}}
\def\fl{\mathsf{fl}}
\def\Flat{\operatorname{Flat}}
\def\Gammam{\Gamma_{\fm}}
\def\Gammamo{\Gamma_{{\fm}^{\opp}}}
\def\Gamman{\Gamma_{\fn}}
\def\Gammano{\Gamma_{{\fn}^{\opp}}}
\def\GInj{\operatorname{GInj}}
\def\GKdim{\operatorname{GKdim}}
\def\GL{\operatorname{GL}}
\def\gldim{\operatorname{gldim}}
\def\GMor{\mathsf{GMor}}
\def\GProj{\mathsf{GProj}}
\def\gr{\mathsf{gr}}
\def\Gr{\operatorname{Gr}}
\def\grade{\operatorname{grade}}
\newcommand{\green}{\color{green}}
\def\GrExt{\operatorname{Ext}}
\def\GrHom{\operatorname{Hom}}
\def\grmod{\mathsf{gr}}
\def\GrMod{\mathsf{Gr}}
\def\GrVect{\operatorname{GrVect}}
\def\h{\operatorname{h}}
\def\H{\operatorname{H}}
\def\hd{\operatorname{hd}}
\def\hdet{\operatorname{hdet}}
\def\HGammam{\operatorname{H}_{\fm}}
\def\HGammamopp{\operatorname{H}_{\fm^{\opp}}}
\def\Hom{\operatorname{Hom}}
\def\id{\operatorname{id}}
\def\Image{\operatorname{Im}}
\def\inc{\operatorname{inc}}
\def\ind{\operatorname{ind}}
\def\indec{\operatorname{ind}}
\def\index{\operatorname{index}}
\def\indexhigher{\operatorname{index}^{\,\tiny{\pentagon}}}
\def\inf{\operatorname{inf}}
\def\infmod{\operatorname{i}}
\def\inj{\operatorname{inj}}
\def\Inj{\operatorname{Inj}}
\def\inv{\operatorname{inv}}
\def\Irr{\operatorname{Irr}}
\def\K{\operatorname{K}}
%\def\K0{\operatorname{K}_0}
\def\Kdim{\operatorname{Kdim}}
\def\Ker{\operatorname{Ker}}
\def\kfd{\operatorname{k.\!fd}}
\def\kid{\operatorname{k.\!id}}
\def\kpd{\operatorname{k.\!pd}}
\def\Ksp{\operatorname{K}_0^{\operatorname{sp}}}
\def\lcd{\operatorname{lcd}}
\def\lAnn{\operatorname{leftann}}
\def\length{\operatorname{length}}
\def\lf{\operatorname{lf}}
\def\lleft{\operatorname{left}}
\newcommand\LTensor[1]{\overset{{\rm L}}{\underset{#1}{\otimes}}}
\def\Mob{\operatorname{M\ddot{o}b}}
\def\mod{\operatorname{mod}}
\def\Mod{\operatorname{Mod}}
\def\Morph{\mathsf{Mor}}
\def\obj{{\operatorname{obj}}}
\def\opp{\operatorname{op}}
\def\ord{\operatorname{ord}}
\def\os{\operatorname{os}}
\def\pcd{\operatorname{pcd}}
\def\pd{\operatorname{pd}}
\def\prj{\operatorname{prj}}
\def\Prj{\operatorname{Prj}}
%\def\prod{\operatorname{prod}}
%\def\Prod{\operatorname{Prod}}
\def\proj{\operatorname{proj}}
\def\Proj{\operatorname{Proj}}
\def\PSL2{\operatorname{PSL}_2}
\def\qgr{\mathsf{qgr}\,}
\def\QGr{\mathsf{QGr}\,}
\def\R{\operatorname{R}\!}
\def\rad{\operatorname{rad}}
\def\rank{\operatorname{rank}}
\def\Res{\operatorname{Res}}
\def\RGammam{\operatorname{R}\!\Gamma_{\fm}}
\def\RGammamo{\operatorname{R}\!\Gamma_{{\fm}^{\opp}}}
\def\RGamman{\operatorname{R}\!\Gamma_{\fn}}
\def\RGammano{\operatorname{R}\!\Gamma_{{\fn}^{\opp}}}
\def\RHom{\operatorname{RHom}}
\def\rk{\mathpzc{k}}
\def\rl{\mathpzc{l}}
\def\RNH{\mathsf{R \mhyphen Nbhd}}
\def\Rpos{\BR_{>0}}
\def\rright{\operatorname{right}}
\def\s{\operatorname{s}}
\def\skewtimes{\ltimes\!}
\def\SL2{\operatorname{SL}_2}
\def\Sp{\operatorname{Sp}}
\def\Spec{\operatorname{Spec}}
\def\split{\operatorname{split}}
\def\stab{\mathsf{stab}}
\def\sup{\operatorname{sup}}
\def\Supp{\operatorname{Supp}}
\newcommand\Tensor[1]{\,{\underset{#1}{\otimes}}\,}
\def\thick{\operatorname{thick}}
\def\Tor{\operatorname{Tor}}
\def\tors{\operatorname{tors}}
\def\Tot{\operatorname{Tot}}
\newcommand\TPath[4]{\cP(#2,#3,#4)}
\def\TP{\mathsf{TorsPair}}
\def\Tr{\operatorname{Tr}}
\def\Uexact{\cU\mbox{\rm -exact}}
\def\Vect{\mathsf{Vect}}
\def\weq{\operatorname{weq}}
\def\Xexact{\cX\mbox{\rm -exact}}
\def\Z{\operatorname{Z}}

%%%%%%%%%%%%%%%%%%%%%%%%%%%%%%%%%%%%%%%%%%%%%%%%%%%%%
%%%%%%%%%%%%%%%%%%%%%%%%%%%%%%%%%%%%%%%%%%%%%%%%%%%%%
%Abbreviations specific to this paper to go here
%%%%%%%%%%%%%%%%%%%%%%%%%%%%%%%%%%%%%%%%%%%%%%%%%%%%%
%%%%%%%%%%%%%%%%%%%%%%%%%%%%%%%%%%%%%%%%%%%%%%%%%%%%%

%Let equation numbers have the form (1.a), (1.b), (2.a), ...
%
\numberwithin{equation}{section}
\renewcommand{\theequation}{\arabic{section}.\arabic{equation}}
%\renewcommand{\theequation}{\arabic{equation}}


%Let subsection numbers have the form (1.a), (1.b), ..., (2.a), (2.b), ...
%
%\renewcommand{\thesubsection}{(\arabic{section}.\roman{subsection})}
\renewcommand{\thesubsection}{\arabic{section}.\arabic{subsection}}
%\renewcommand{\thesubsection}{\Alph{subsection}}


%Use Roman numbering in enumerate environment
%
\renewcommand{\labelenumi}{(\roman{enumi})}


%The following defines environments for Theorems, Lemmas, etc.
%
\newtheorem{Lemma}{Lemma}[section]
%For simple numbering of Theorems, Lemmas etc., uncomment the
%following line
%\renewcommand{\theLemma}{\arabic{Lemma}}
\newtheorem{Theorem}[Lemma]{Theorem}
\newtheorem{Proposition}[Lemma]{Proposition}
\newtheorem{Corollary}[Lemma]{Corollary}

\theoremstyle{definition}
\newtheorem{Definition}[Lemma]{Definition}
\newtheorem{Setup}[Lemma]{Setup}
\newtheorem{Conjecture}[Lemma]{Conjecture}
\newtheorem{Condition}[Lemma]{Condition}
\newtheorem{Conditions}[Lemma]{Conditions}
\newtheorem{Construction}[Lemma]{Construction}
\newtheorem{Remark}[Lemma]{Remark}
\newtheorem{Remarks}[Lemma]{Remarks}
\newtheorem{Observation}[Lemma]{Observation}
\newtheorem{Example}[Lemma]{Example}
\newtheorem{Convention}[Lemma]{Convention}
\newtheorem{Notation}[Lemma]{Notation}
\newtheorem{Question}[Lemma]{Question}
\newtheorem{Empty}[Lemma]{}

%The following lines make special environments
%for (e.g.) Theorems A, B, C,... in the introduction, 
%invoked with (e.g.) \begin{ThmIntro} ... \end{ThmIntro}.

\theoremstyle{theorem}
\newtheorem{ThmIntro}{Theorem}
\newtheorem{CorIntro}[ThmIntro]{Corollary}
\theoremstyle{definition}
\newtheorem{DefIntro}[ThmIntro]{Definition}
\renewcommand{\theDefIntro}{\Alph{DefIntro}}
\renewcommand{\theThmIntro}{\Alph{ThmIntro}}
\renewcommand{\theCorIntro}{\Alph{CorIntro}}

%\theoremstyle{bfupright head,upright body}
\newtheorem{bfhpg}[Lemma]{}               
\newtheorem*{bfhpg*}{}

\newtheorem*{ClThm}{Classic Theorem}

%The following lines make a new type of theorem, altbfhpg, used in the document.
\usepackage{amsthm}
\usepackage{thmtools}
\declaretheoremstyle[
notefont=\bfseries, notebraces={}{},
bodyfont=\normalfont,
%headformat=\NUMBER~\NAME:\NOTE
headformat=\NUMBER~\NOTE,
%headpunct={:}
headpunct={}
]{foobar}
\declaretheorem[style=foobar,numberlike=Lemma]{altbfhpg}


% Useful description list that indents based on the length of a
% parameter, presumably the longest label.
\newenvironment{VarDescription}[1]%
  {\begin{list}{}{\renewcommand{\makelabel}[1]{\hfill\text{##1}}%
    \settowidth{\labelwidth}{\textbf{#1:}}%
    \setlength{\leftmargin}{\labelwidth}\addtolength{\leftmargin}{\labelsep}}}%
  {\end{list}}


%The following makes 2020 available as a parameter for the \subjclass command.
\makeatletter
\@namedef{subjclassname@2020}{%
  \textup{2020} Mathematics Subject Classification}
\makeatother




\begin{document}

\setlength{\parindent}{0pt}
\setlength{\parskip}{7pt}


\title[Splitting higher torsion classes]{Splitting torsion classes in higher homological algebra}


\author{Erlend D.\ B\o rve}

\address{Department of Mathematics, Aarhus University, Ny Munkegade 118, 8000 Aarhus C, Denmark}
\email{erlend.d.borve@math.au.dk}

%\urladdr{https://bervinator.github.io}


\author{Peter J\o rgensen}

\address{Department of Mathematics, Aarhus University, Ny Munkegade 118, 8000 Aarhus C, Denmark}
\email{peter.jorgensen@math.au.dk}

\urladdr{https://sites.google.com/view/peterjorgensen}


\author{Mads Hustad Sand\o y}

\address{Department of Mathematics, Aarhus University, Ny Munkegade 118, 8000 Aarhus C, Denmark}
\email{mads.sandoy@ntnu.no}

\urladdr{https://sites.google.com/view/mads-h-sandoy-math}


%\thanks{Date: \today. A thank you would go here}

\keywords{$d$-abelian category, $d$-Auslander--Reiten theory, $d$-torsion class, higher Auslander algebra of Dynkin type $A$, lattice, spine}

\subjclass[2020]{05E10, 06A07, 16S90, 18E99}

%05B45: Combinatorial aspects of tessellation and tiling problems
%05E10: Combinatorial aspects of representation theory
%05E16: Combinatorial aspects of groups and algebras
%05E40: Combinatorial aspects of commutative algebra
%05E45: Combinatorial aspects of simplicial complexes
%05E99: Algebraic combinatorics / None of the above, but in this section
%06A07: Combinatorics of partially ordered sets
%13F60: Cluster algebras
%16E10: Homological dimension in associative algebras
%16E35: Derived categories and associative algebras
%16E45: Differential graded algebras and applications (associative algebraic aspects)
%16G10: Representations of associative Artinian rings 
%16G60: Representation type (finite, tame, wild, etc.) of associative algebras
%16G70: Auslander-Reiten sequences (almost split sequences) and Auslander-Reiten quivers
%16S90: Torsion theories; radicals on module categories (associative algebraic aspects) {For radicals of rings, see 16Nxx}
%18A20: Epimorphisms, monomorphisms, special classes of morphisms, null morphisms
%18E10: Abelian categories, Grothendieck categories
%18E35: Localization of categories, calculus of fractions
%18E40: Torsion theories, radicals
%18E99: Categorical algebra / None of the above, but in this section
%18G05: Projectives and injectives (category-theoretic aspects) [See also 13C10, 13C11, 16D40, 16D50]
%18G35: Chain complexes (category-theoretic aspects), dg categories [See also 14F07, 18G80, 55U15]
%18G80: Derived categories, triangulated categories
%18G99: Homological algebra in category theory, derived categories and functors / None of the above, but in this section 
%18N40: Homotopical algebra, Quillen model categories, derivators
%51M20: Polyhedra and polytopes; regular figures, division of spaces
%55P62: Rational homotopy theory



\begin{abstract} 

Let $d \geqslant 1$ be an integer, $\cM$ a suitable $d$-abelian category.  There is a notion of higher torsion classes in $\cM\!$, also known as $d$-torsion classes.  

\medskip
\noindent
We generalise a classic theorem of Hoshino by providing criteria for a $d$-torsion class $\cU$ to be splitting,  notably that $\cU$ is stable under $\tau_d^-$, the inverse $d$-Auslander--Reiten translation.

\medskip
\noindent
We apply the criteria to show that under additional assumptions, which are satisfied in Dynkin type $A$, the elements of the spine of the lattice of $d$-torsion classes are precisely the splitting $d$-torsion classes.  The spine consists of the elements which belong to a chain of maximal length. 

\end{abstract}




\maketitle




\setcounter{section}{-1}
\section{Introduction}
\label{sec:introduction}


Higher homological algebra was founded by Iyama, see \cite{Iyama_Auslander-correspondence}, \cite{Iyama_Cluster-tilting-for-higher-Auslander-algebras}, \cite{Iyama_Higher-dimensional-AR-theory}.  A subsequent key step was taken by Jasso who introduced $d$-abelian categories for integers $d \geqslant 1$, see  \cite[def.\ 3.1]{Jasso_n-abelian}.  A $1$-abelian category is just an abelian category in the classic sense, while a $d$-abelian category has complexes of length $d$ instead of kernels and cokernels.

Higher torsion classes in $d$-abelian categories generalise torsion classes in classic abelian categories.  They are also known as $d$-torsion classes and were introduced in \cite[def.\ 1.1]{Jorgensen_Torsion-classes-in-higher-homological-algebra}, where splitting $d$-torsion classes were defined too.  See Remark \ref{rmk:d-torsion_classes}.  

Our first main result is the following generalisation of the classic characterisation of splitting torsion classes due to Hoshino; see \cite[prop.\ 1]{Hoshino_On-splitting-torsion-theories-induced-by-tilting-modules} and \cite[prop.\ VI.1.7]{Assem-Simson-Skowronski_Elements-I}.  Let $\cM$ be a $d$-abelian category arising as a $d$-cluster tilting subcategory of $\mod\,A$ for a finite dimensional algebra $A$ over an algebraically closed field $k$; see \cite[def.\ 2.2]{Iyama_Higher-dimensional-AR-theory} and \cite[def.\ 3.14]{Jasso_n-abelian}.  


\begin{ThmIntro}
[= Theorem \ref{thm:generalisation_of_ASS_VI1.7}]
\label{thm:A}
Let $\cU \subseteq \cM$ be a $d$-torsion class.  The following two conditions are equivalent:
\begin{enumerate}
\setlength\itemsep{4pt}

  \item  $\cU$ is splitting.
  
  \item  $\tau_d^-\cU \subseteq \cU\!$, where $\tau_d^-$ is the inverse $d$-Auslander--Reiten translation. 

\end{enumerate}
Now assume $d \geqslant 2$.  Then conditions (i) and (ii) are also equivalent to the following:
\begin{enumerate}
\setcounter{enumi}{2}
\setlength\itemsep{4pt}

  \item  If $\delta = 0 \xrightarrow{} U \xrightarrow{} M \xrightarrow{} V^1 \xrightarrow{} \cdots \xrightarrow{} V^d \xrightarrow{} 0$ is a canonical $d$-exact sequence with respect to $\cU\!$, then
\[  
  \overline{ \Hom }_A( \widetilde{ U },\tau_dV^{ d-1 } ) \xrightarrow{} \overline{ \Hom }_A( \widetilde{ U },\tau_dV^{ d } )
\]  
is surjective for each $\widetilde{ U } \in \cU\!$, where $\tau_d$ is the $d$-Auslander--Reiten translation and $\overline{ \Hom }_A$ is injectively stabilised $\Hom_A$.  

\end{enumerate}
\end{ThmIntro}

Note that the classic characterisation of splitting torsion classes also provides a condition on the corresponding torsion free classes.  This is not possible for $d$-torsion classes, which do not in general have corresponding $d$-torsion free classes.  However, a substitute is provided by condition (iii) in Theorem \ref{thm:A}.

Our second main result concerns the partially ordered set of $d$-torsion classes.  It was proved by August, Haugland, Jacobsen, Kvamme, Palu, and Treffinger that it is a lattice; see \cite[thm.\ 4.3]{August-et-al_A-characterisation-of-higher-torsion-classes}.  The spine of a lattice was originally defined for trim lattices by Thomas in \cite[p.\ 257]{Thomas_An-analogue-of-distributivity-for-ungraded-lattices}; it consists of the elements which belong to a chain of maximal length.  The spine of the lattice of classic torsion classes has recently been investigated by Asai, Iyama, Mousavand, and Paquette; see \cite[thm.\ 1.7]{Asai-Iyama-Mousavand-Paquette_Brick-splitting-torsion-pairs-and-trim-lattices}.  We show the following result in the higher case.

\begin{ThmIntro}
[= Theorem \ref{thm:spine}]
\label{thm:B}
Assume that $\cM$ has no cycles in the sense of Remark \ref{rmk:cycles} and that it has only finitely many isomorphism classes of indecomposable modules.  

Then the elements of the spine of the lattice of $d$-torsion classes of $\cM$ are precisely the splitting $d$-torsion classes.
\end{ThmIntro}

Theorem \ref{thm:B} applies in particular to Dynkin type $A$; see Example \ref{exa:Type_A}.

We refer to \cite{Asadollahi-Jorgensen-Schroll-Treffinger_On-higher-torsion-classes}, \cite{August-et-al_A-characterisation-of-higher-torsion-classes}, \cite{August-et-al_Higher-torsion-classes-tau-d-tilting-theory-and-silting-complexes}, \cite{Fedele-Jorgensen-Shah_The-index-in-d-exact-categories}, \cite{Kvamme_Higher-extension-closure-and-d-exact-categories}, \cite{Segovia_P-comma-phi-Tamari-and-higher-torsion-lattices-of-type-A}, \cite{Segovia_P-comma-phi-Tamari-lattices} for other recent contributions to the theory of $d$-torsion classes.

We end the introduction with a blanket setup and a few background items.


\begin{Setup}
The following will be fixed throughout.
\begin{enumerate}
\setlength\itemsep{4pt}

  \item  $k$ is an algebraically closed field and $d \geqslant 1$ is an integer.
  
  \item  $A$ is a finite dimensional $k$-algebra and $\mod A$ is the abelian category of finite dimensional right $A$-modules.
  
  \item  $\dual( - ) = \Hom_k( -,k )$ is the $k$-linear duality functor.

  \item  $\cM \subseteq \mod\,A$ is a $d$-cluster tilting subcategory; see \cite[def.\ 2.2]{Iyama_Higher-dimensional-AR-theory} and \cite[def.\ 3.14]{Jasso_n-abelian}.   
  
  \item  If $\cX \subseteq \cM$ is an additive subcategory, then we write
\[
  \cX^{ \perp }
  = \{\, Y \in \cM \,|\, \Hom_A( \cX,Y ) = 0 \,\}.
\]

\end{enumerate}
\end{Setup}


\begin{Remark}
[$d$-abelian categories]
Under the setup, $\cM$ is a $d$-abelian category by \cite[thm.\ 3.16]{Jasso_n-abelian}.  The notion of $d$-exact sequences in $\cM$ will be used repeatedly; see \cite[def.\ 2.4]{Jasso_n-abelian}.  It follows from \cite[props.\ 2.2 and 2.6]{Jasso-Kvamme_Introduction-to_higher_AR-theory} that a $d$-exact sequence in $\cM$ is simply an exact sequence in $\mod\,A$,
\[
  0
  \xrightarrow{} M^0
  \xrightarrow{} M^1
  \xrightarrow{} \cdots
  \xrightarrow{} M^d
  \xrightarrow{} M^{ d+1 }
  \xrightarrow{} 0,
\]
with $M^i \in \cM$ for each $i$.
\end{Remark}


\begin{Remark}
[cycles]
\label{rmk:cycles}
Recall from \cite[p.\ 313]{Auslander-Reiten-Smalo_Representation-theory-of-Artin-algebras} that a cycle from $M$ to $M$ in $\cM$ is a sequence $M \xrightarrow{ \mu_0 } M_1 \xrightarrow{ \mu_1 } \cdots \xrightarrow{ \mu_{ t-2 } } M_{ t-1 } \xrightarrow{ \mu_{ t-1 } } M$ with $t \geqslant 1$, where $M$ and each $M_i$ is indecomposable in $\cM\!$, and each $\mu_i$ is not zero and not an isomorphism.  Note that if $t = 1$, then the cycle has the form $M \xrightarrow{ \mu_0 } M$.
\end{Remark}


\begin{Remark}
[$d$-torsion classes]
\label{rmk:d-torsion_classes}
${\;}$
\begin{enumerate}
\setlength\itemsep{4pt}

  \item  Let $\cU \subseteq \cM$ be an additive subcategory.  Let $M \in \cM$ be given.  A canonical $d$-exact sequence of $M$ with respect to $\cU$ is a $d$-exact sequence 
\begin{equation}
\label{equ:canonical_sequence}
  0
  \xrightarrow{} U
  \xrightarrow{ \upsilon } M
  \xrightarrow{} V^1
  \xrightarrow{} \cdots
  \xrightarrow{} V^d
  \xrightarrow{} 0
\end{equation}
in $\cM$ with $U \in \cU$ such that
\[
  0
  \xrightarrow{} \Hom_A( \widetilde{ U },V^1 )
  \xrightarrow{} \cdots 
  \xrightarrow{} \Hom_A( \widetilde{ U },V^d )
  \xrightarrow{} 0
\]
is exact for each $\widetilde{ U } \in \cU\!$.  In such a sequence, $M$ determines $\upsilon$ up to isomorphism since $\upsilon$ is a $\cU$-cover of $M$; see \cite[lem.\ 2.7(i)]{Jorgensen_Torsion-classes-in-higher-homological-algebra}.

  \item  A $d$-torsion class $\cU \subseteq \cM$ is an additive subcategory such that each $M \in \cM$ has a canonical $d$-exact sequence with respect to $\cU\!$.
  
  \item  A $d$-torsion class $\cU \subseteq \cM$ is called splitting if each $M \in \cM$ has a canonical $d$-exact sequence \eqref{equ:canonical_sequence} with respect to $\cU$ where $\upsilon$ is a split monomorphism.  In this case, every canonical $d$-exact sequence with respect to $\cU$ has $\upsilon$ a split monomorphism because $M$ determines $\upsilon$ up to isomorphism.

  \item  The set of $d$-torsion classes in $\cM$ is denoted $\dtors\,\cM$.  It is clearly a partially ordered set under inclusion.  This partially ordered set is, in fact, a lattice; that is, it has meets and joins.  This holds by \cite[thm.\ 4.3]{August-et-al_A-characterisation-of-higher-torsion-classes}, which also states that meets are given by intersection.

\end{enumerate}
\end{Remark}


\begin{Remark}
[$d$-Auslander--Reiten theory]
\label{rmk:tau_d}
${\;}$
\begin{enumerate}
\setlength\itemsep{4pt}

  \item  The projective stable module category $\underline{ \mod }\,A$ is the na\"{i}ve quotient of $\mod A$ by the ideal of homomorphisms factoring through a projective module, and its $\Hom$ functor is $\underline{ \Hom }_A$.  There is a full subcategory $\underline{ \cM } \subseteq \underline{ \mod }\,A$ induced by $\cM \subseteq \mod\,A$.
  
  \item  The injective stable module category $\overline{ \mod }\,A$ is the na\"{i}ve quotient of $\mod A$ by the ideal of homomorphisms factoring through an injective module, and its $\Hom$ functor is $\overline{ \Hom }_A$.  There is a full subcategory $\overline{ \cM } \subseteq \overline{ \mod }\,A$ induced by $\cM \subseteq \mod\,A$.

  \item  The $d$-Auslander--Reiten translation $\tau_d$ and its left adjoint $\tau_d^-$ are functors
\[
  \begin{tikzcd}
	{\underline{ \mod }\,A} & {\overline{ \mod }\,A}\lefteqn{,}
	\arrow["{\tau_d}"', shift right=1, from=1-1, to=1-2]
	\arrow["{\tau_d^-}"', shift right=1, from=1-2, to=1-1]
  \end{tikzcd}
\]  
see \cite[thm.\ 2.3.1(1)]{Iyama_Auslander-correspondence} and \cite[thm.\ 1.4.1]{Iyama_Higher-dimensional-AR-theory}.  They restrict to pseudo-inverse functors
\[
  \begin{tikzcd}
	{\underline{ \cM }} & {\overline{ \cM }\!.}
	\arrow["{\tau_d}"', shift right=1, from=1-1, to=1-2]
	\arrow["{\tau_d^-}"', shift right=1, from=1-2, to=1-1]
  \end{tikzcd}
\]  

  \item  A $d$-exact sequence $\delta$ in $\cM$ gives rise to co- and contravariant defect functors $\delta_*$ and $\delta^*$ defined in \cite[def.\ 3.1]{Jasso-Kvamme_Introduction-to_higher_AR-theory}.  By \cite[thm.\ 3.8]{Jasso-Kvamme_Introduction-to_higher_AR-theory} they satisfy the higher defect formula
\begin{equation}
\label{equ:defect_formula}
  \dual\!\delta^*( - ) \cong \delta_*\big( \tau_d( - ) \big),
\end{equation}
where both sides are functors defined on $\underline{ \cM }$, and its dual
\begin{equation}
\label{equ:dual_defect_formula}
  \dual\!\delta_*( - ) \cong \delta^*\big( \tau^-_d( - ) \big),
\end{equation}
where both sides are functors defined on $\overline{ \cM }\!$.

  \item  If $M \in \cM$ is given, then $\tau_d M$ is only defined up to isomorphism in $\overline{ \cM }$, but we make it defined up to isomorphism in $\cM$ by picking it not to have any non-zero injective summands.  Similarly, $\tau_d^- M$ is only defined up to isomorphism in $\underline{ \cM }$, but we make it defined up to isomorphism in $\cM$ by picking it not to have any non-zero projective summands.  

\medskip
\noindent
With this stipulation, $\tau_d$ and $\tau_d^-$ induce bijections between the isomorphism classes of non-projective indecomposable modules in $\cM$ and the isomorphism classes of non-injective indecomposable modules in $\cM\!$, see \cite[thm.\ 2.3]{Iyama_Higher-dimensional-AR-theory}.

\end{enumerate}
\end{Remark}




\section{Homological characterisation of splitting $d$-torsion classes}
\label{sec:splitting}


\begin{Lemma}
\label{lem:criterion_for_splitting}
Let $\cU \subseteq \cM$ be an additive subcategory.  The following conditions are equivalent:
\begin{enumerate}
\setlength\itemsep{4pt}

  \item  $\cU$ is a splitting $d$-torsion class.

  \item  Each $M \in \cM$ can be written $M = U \coprod V$ with $U \in \cU$ and $V \in \cU^{ \perp }$.

  \item  Each indecomposable $M \in \cM$ is either in $\cU$ or in $\cU^{ \perp }$.
  
\end{enumerate}
\end{Lemma}

\begin{proof}
(i)$\Rightarrow$(ii):  Assume that (i) is satisfied.  Let $M \in \cM$ be given.  Pick a canonical $d$-exact sequence \eqref{equ:canonical_sequence} of $M$ with respect to $\cU$ where $\upsilon$ is a split monomorphism.  Up to isomorphism, we can write $M = U \coprod V$ and view $\upsilon$ as the coproduct inclusion $U \xrightarrow{} M$.  Since $\upsilon$ is a $\cU$-cover by \cite[lem.\ 2.7(i)]{Jorgensen_Torsion-classes-in-higher-homological-algebra}, this implies $V \in \cU^{ \perp }$ which proves (ii).

(ii)$\Rightarrow$(i):  Assume that (ii) is satisfied.  Let $M \in \cM$ be given.  Write $M = U \coprod V$ with $U \in \cU$ and $V \in \cU^{ \perp }$.  Let $U \xrightarrow{ \iota } M$ be the coproduct inclusion.  There is a $d$-exact sequence 
\begin{equation}
\label{equ:lem:criterion_for_splitting:20}
  0
  \xrightarrow{} U
  \xrightarrow{ \iota } M
  \xrightarrow{} V
  \xrightarrow{} 0
  \xrightarrow{} \cdots
  \xrightarrow{} 0
\end{equation}
in $\cM$ by \cite[prop.\ 2.2]{Jasso-Kvamme_Introduction-to_higher_AR-theory} which can be used as the canonical $d$-exact sequence \eqref{equ:canonical_sequence}.  Since $\iota$ is a split monomorphism, this proves (i).

(ii)$\Leftrightarrow$(iii) is immediate because $\cM$ is a Krull--Schmidt category.
\end{proof}


\begin{Lemma}
\label{lem:criterion_for_contractible}
Let $\cB$ be an additive category and let $\K( \cB )$ be the homotopy category of chain complexes over $\cB\!$.  Let $B$ be a bounded chain complex over $\cB\!$.  The following conditions are equivalent:
\begin{enumerate}
\setlength\itemsep{4pt}

  \item  $B$ is isomorphic to zero in $\K( \cB )$.

  \item  For each $\widetilde{ B } \in \cB$ the complex $\Hom_{ \cB }( \widetilde{ B },B )$ is exact.
  
  \item  For each $\widetilde{ B } \in \cB$ the complex $\Hom_{ \cB }( B,\widetilde{ B } )$ is exact.

\end{enumerate}
\end{Lemma}

\begin{proof}
If $i \in \BZ$ then \cite[thm.\ 2.3.10]{Christensen-Foxby-Holm_Derived-category-methods} implies $\H_i\!\Hom_{ \cB }( \widetilde{ B },B ) \cong \Hom_{ \K( \cB ) }( \Sigma^i\widetilde{ B },B )$, where $\Sigma$ is the suspension functor of $\K( \cB )$ and $\widetilde{ B }$ is viewed as a complex concentrated in degree zero.  Hence (ii) is equivalent to the following condition.
\begin{itemize}
\setlength\itemsep{4pt}

  \item[(ii')]  For each $i \in \BZ$ and $\widetilde{ B } \in \cB$ we have $\Hom_{ \K( \cB ) }( \Sigma^i\widetilde{ B },B ) = 0$.
  
\end{itemize}
To prove (i)$\Leftrightarrow$(ii) it is hence enough to prove (i)$\Leftrightarrow$(ii'), and the implication (i)$\Rightarrow$(ii') is clear.  To prove (ii')$\Rightarrow$(i), note that in the triangulated category $\K( \cB )$, the bounded complex $B$ can be built from objects of the form $\Sigma^i\widetilde{ B }$ using finitely many hard truncation triangles.  Hence (ii') implies $\Hom_{ \K( \cB ) }( B,B ) = 0$, which is equivalent to (i).

The equivalence (i)$\Leftrightarrow$(iii) is proved similarly.
\end{proof}


\begin{Lemma}
\label{lem:splitting_implies_Ext^d_zero}
Let $\cU \subseteq \cM$ be a splitting $d$-torsion class.  Let $U \in \cU$ and $V \in \cU^{ \perp }$ be given.  Then each $d$-exact sequence in $\cM$ of the form
\begin{equation}
\label{equ:lem:splitting_implies_Ext^d_zero:10}
  0
  \xrightarrow{} U
  \xrightarrow{} M^1
  \xrightarrow{} \cdots
  \xrightarrow{} M^d
  \xrightarrow{} V
  \xrightarrow{} 0
\end{equation}
is split exact.
\end{Lemma}

\begin{proof}
By Lemma \ref{lem:criterion_for_splitting}(ii) we can write $M^i = U^i \oplus V^i$ with $U^i \in \cU$ and $V^i \in \cU^{ \perp }$ for each $i$.  Hence \eqref{equ:lem:splitting_implies_Ext^d_zero:10} can be spun into the following commutative diagram where the vertical morphisms are inclusions and projections.
\[
  \begin{tikzcd}
	0 & U & {U^1} & \cdots & {U^d} & 0 & 0 \\
	0 & U & {U^1 \coprod V^1} & \cdots & {U^d \coprod V^d} & V & 0 \\
	0 & 0 & {V^1} & \cdots & {V^d} & V & 0
	\arrow[from=1-1, to=1-2]
	\arrow[from=1-2, to=1-3]
	\arrow[from=1-2, to=2-2]
	\arrow[from=1-3, to=1-4]
	\arrow[from=1-3, to=2-3]
	\arrow[from=1-4, to=1-5]
	\arrow[from=1-5, to=1-6]
	\arrow[from=1-5, to=2-5]
	\arrow[from=1-6, to=1-7]
	\arrow[from=1-6, to=2-6]
	\arrow[from=2-1, to=2-2]
	\arrow[from=2-2, to=2-3]
	\arrow[from=2-2, to=3-2]
	\arrow[from=2-3, to=2-4]
	\arrow[from=2-3, to=3-3]
	\arrow[from=2-4, to=2-5]
	\arrow[from=2-5, to=2-6]
	\arrow[from=2-5, to=3-5]
	\arrow[from=2-6, to=2-7]
	\arrow[from=2-6, to=3-6]
	\arrow[from=3-1, to=3-2]
	\arrow[from=3-2, to=3-3]
	\arrow[from=3-3, to=3-4]
	\arrow[from=3-4, to=3-5]
	\arrow[from=3-5, to=3-6]
	\arrow[from=3-6, to=3-7]
  \end{tikzcd}
\]
Extending by zeroes to the left and right, each line in the diagram is a chain complex over $\cM\!$.  The whole diagram (rotated anticlockwise by $90$ degrees) is a semisplit short exact sequence
\begin{equation}
\label{equ:lem:splitting_implies_Ext^d_zero:20}
  0 \xrightarrow{} U^* \xrightarrow{} E^* \xrightarrow{} V^* \xrightarrow{} 0
\end{equation}
of chain complexes over $\cM\!$.

Let $\widetilde{ U } \in \cU$ be given.  Then the complex
\begin{equation}
\label{equ:lem:splitting_implies_Ext^d_zero:30}
  \Hom_A( \widetilde{ U },V^* )
\end{equation}
is exact; in fact, it is zero since each object in $V^*$ is in $\cU^{ \perp }$.  The complex
\begin{equation}
\label{equ:lem:splitting_implies_Ext^d_zero:40}
  \Hom_A( \widetilde{ U },E^* )
\end{equation}
is also exact because the sequence
\[
  0
  \xrightarrow{} \Hom_A( \widetilde{ U },U )
  \xrightarrow{} \Hom_A( \widetilde{ U },U^1 \coprod V^1 )
  \xrightarrow{} \cdots
  \xrightarrow{} \Hom_A( \widetilde{ U },U^d \coprod V^d )
  \xrightarrow{} \Hom_A( \widetilde{ U },V )  
\]
is exact since $E^*$ is a $d$-exact sequence, and because we have $\Hom_A( \widetilde{ U },V ) = 0$ by $\widetilde{ U } \in \cU$ and $V \in \cU^{ \perp }$.  Since the complexes \eqref{equ:lem:splitting_implies_Ext^d_zero:30} and \eqref{equ:lem:splitting_implies_Ext^d_zero:40} are exact, the long exact sequence induced by \eqref{equ:lem:splitting_implies_Ext^d_zero:20} implies that $\Hom_A( \widetilde{ U },U^* )$ is exact.  Since this holds for each $\widetilde{ U } \in \cU$ while each object in $U^*$ is in $\cU$, Lemma \ref{lem:criterion_for_contractible} implies that $U^*$ is isomorphic to zero in $\K( \cU )$, hence isomorphic to zero in $\K( \mod A )$.

A symmetric argument using $\Hom_A( -,\widetilde{ V } )$ with $\widetilde{ V } \in \cU^{ \perp }$ shows $V^*$ isomorphic to zero in $K( \mod A )$.

Since \eqref{equ:lem:splitting_implies_Ext^d_zero:20} is a semisplit short exact sequence of chain complexes over $\cM$ and hence over $\mod A$, it induces a triangle $U^* \xrightarrow{} E^* \xrightarrow{} V^* \xrightarrow{} \Sigma U^*$ in the triangulated category $\K( \mod A )$.  Hence $E^*$ is isomorphic to zero in $\K( \mod A )$ because so are $U^*$ and $V^*$.  But $E^*$ is \eqref{equ:lem:splitting_implies_Ext^d_zero:10} extended by zeroes to the left and right, so \eqref{equ:lem:splitting_implies_Ext^d_zero:10} is split exact.
\end{proof}


\begin{Lemma}
\label{lem:splitting_vs_defects}
Let $\cU \subseteq \cM$ be a $d$-torsion class.  The following conditions are equivalent:
\begin{enumerate}
\setlength\itemsep{4pt}

  \item  $\cU$ is splitting.
  
  \item  If $\delta$ is a canonical $d$-exact sequence with respect to $\cU$ then $\delta_*( \cU ) = 0$.

  \item  If $\delta$ is a canonical $d$-exact sequence with respect to $\cU$ then $\delta^*( \tau_d^-\cU ) = 0$.

\end{enumerate}
\end{Lemma}

\begin{proof}
We will use the notation
\[
  \delta 
  \;\;=\;\; 0
  \xrightarrow{} U
  \xrightarrow{ \upsilon } M
  \xrightarrow{} V^1
  \xrightarrow{} \cdots
  \xrightarrow{} V^d
  \xrightarrow{} 0.
\]
By \cite[def.\ 3.1]{Jasso-Kvamme_Introduction-to_higher_AR-theory} the covariant defect functor $\delta_*$ is defined by the exact sequence
\[
  \Hom_A( M,- )
  \xrightarrow{ \upsilon^* }
  \Hom_A( U,- )
  \xrightarrow{}
  \delta_*( - )
  \xrightarrow{}
  0.
\]

(i)$\Rightarrow$(ii):  Assume that (i) is satisfied.  For each canonical $d$-exact sequence $\delta$, the homomorphism $\upsilon$ is a split monomorphism whence $\upsilon^*$ is surjective, so we have $\delta_* = 0$ and in particular $\delta_*( \cU ) = 0$.  This proves (ii).

(ii)$\Rightarrow$(i):  Assume that (ii) is satisfied.   For each canonical $d$-exact sequence $\delta$, we have $\delta_*( U ) = 0$ in particular, whence
\[
  \Hom_A( M,U )
  \xrightarrow{ \Hom_A( \upsilon,U ) }
  \Hom_A( U,U )
\]  
is surjective so $\upsilon$ is a split monomorphism.  This proves (i).

(ii)$\Leftrightarrow$(iii) holds by the dual defect formula, Equation \eqref{equ:dual_defect_formula}.
\end{proof}


\begin{Lemma}
\label{lem:surjective_simultaneously}
Let $Y \stackrel{ \upsilon }{ \twoheadrightarrow } Y''$ be a surjection in $\mod\,A$.  For each $X \in \mod\,A$, the induced homomorphisms
\[
  \Hom_A( X,Y ) \xrightarrow{ \alpha } \Hom_A( X,Y'' )
\]
and
\[
  \underline{ \Hom }_A( X,Y ) \xrightarrow{ \beta } \underline{ \Hom }_A( X,Y'' )
\]
are surjective simultaneously.
\end{Lemma}

\begin{proof}
If $\xi$ is a homomorphism in $\Hom_A$, then $\underline{ \xi }$ will denote its class in $\underline{ \Hom }_A$.  

Suppose $\alpha$ is surjective.  Let $\underline{ \xi }'' \in \underline{ \Hom }_A( X,Y'' )$ be given.  Pick $\xi \in \Hom_A( X,Y )$ such that $\alpha( \xi ) = \xi''$.  Then $\beta( \underline{ \xi } ) = \underline{ \xi }''$.  This proves $\beta$ is surjective.

Suppose $\beta$ is surjective.  Let $\xi'' \in \Hom_A( X,Y'' )$ be given.  Pick $\underline{ \xi } \in \underline{ \Hom }_A( X,Y )$ such that $\beta( \underline{ \xi } ) = \underline{ \xi }''$.  This equation means that $\alpha( \xi )$ and $\xi''$ become equal when projected to $\underline{ \Hom }_A( X,Y'' )$, whence $\alpha( \xi ) - \xi''$ factorises through a projective module $P$ as follows.
\[
\begin{tikzcd}
	&& Y \\
	X & P \\
	&& {Y''}
	\arrow["\upsilon", two heads, from=1-3, to=3-3]
	\arrow["\varphi", dashed, from=2-1, to=1-3]
	\arrow[from=2-1, to=2-2]
	\arrow["{\alpha( \xi ) - \xi''}"', from=2-1, to=3-3]
	\arrow[dotted, from=2-2, to=1-3]
	\arrow[from=2-2, to=3-3]
\end{tikzcd}
\]
The dotted arrow exists because $P$ is projective and $\upsilon$ is surjective.  It induces the dashed arrow $\varphi$ whence $\alpha( \xi ) - \xi'' = \upsilon\varphi$.  Hence we have $\xi'' = \alpha( \xi ) - \upsilon\varphi = \alpha( \xi - \varphi )$.  This proves $\alpha$ is surjective.
\end{proof}


\begin{Theorem}
\label{thm:generalisation_of_ASS_VI1.7}
Let $\cU \subseteq \cM$ be a $d$-torsion class.  The following two conditions are equivalent:
\begin{enumerate}
\setlength\itemsep{4pt}

  \item  $\cU$ is splitting.
  
  \item  $\tau_d^-\cU \subseteq \cU\!$.

\end{enumerate}
Now assume $d \geqslant 2$.  Then conditions (i) and (ii) are also equivalent to the following:
\begin{enumerate}
\setcounter{enumi}{2}
\setlength\itemsep{4pt}

  \item  If $\delta = 0 \xrightarrow{} U \xrightarrow{} M \xrightarrow{} V^1 \xrightarrow{} \cdots \xrightarrow{} V^d \xrightarrow{} 0$ is a canonical $d$-exact sequence with respect to $\cU$, then
\[  
  \overline{ \Hom }_A( \widetilde{ U },\tau_dV^{ d-1 } ) \xrightarrow{} \overline{ \Hom }_A( \widetilde{ U },\tau_dV^{ d } )
\]  
is surjective for each $\widetilde{ U } \in \cU$.  

\end{enumerate}
\end{Theorem}

\begin{proof}
If $d=1$ then the result (i)$\Leftrightarrow$(ii) is contained in \cite[prop.\ VI.1.7]{Assem-Simson-Skowronski_Elements-I}, so we assume $d \geqslant 2$ for the rest of the proof.

(i)$\Rightarrow$(ii):  Assume that (i) is satisfied.  Let $U \in \cU$ be indecomposable.  It is enough to show $\tau_d^-U \in \cU$.  If $U$ is injective then $\tau_d^-U = 0$ is clearly in $\cU\!$, so let us assume that $U$ is non-injective.  By \cite[thm.\ 3.3.1]{Iyama_Higher-dimensional-AR-theory} there is a $d$-Auslander--Reiten sequence in $\cM\!$,
\begin{equation}
\label{equ:thm:generalisation_of_ASS_VI1.7:10}
  0
  \xrightarrow{} U 
  \xrightarrow{} M^1 
  \xrightarrow{} \cdots
  \xrightarrow{} M^d
  \xrightarrow{} \tau_d^-U 
  \xrightarrow{} 0.
\end{equation}
Lemma \ref{lem:criterion_for_splitting}(iii) gives $\tau_d^-U \in \cU$ or $\tau_d^-U \in \cU^{ \perp }$.  But the latter would imply that \eqref{equ:thm:generalisation_of_ASS_VI1.7:10} were split exact by Lemma \ref{lem:splitting_implies_Ext^d_zero}.  This is false by definition of $d$-Auslander--Reiten sequences, see \cite[p.\ 33]{Iyama_Higher-dimensional-AR-theory}, so we conclude $\tau_d^-U \in \cU$.  We have proved (ii).

Before continuing, we claim that condition (i) in the theorem is equivalent to the following:
\begin{itemize}
\setlength\itemsep{4pt}

  \item[(i')]  If $\delta = 0 \xrightarrow{} U \xrightarrow{} M \xrightarrow{} V^1 \xrightarrow{} \cdots \xrightarrow{} V^d \xrightarrow{} 0$ is a canonical $d$-exact sequence with respect to $\cU$, then
\[  
  \Hom_A( \tau_d^-\widetilde{ U },V^{ d-1 } )
  \xrightarrow{}
  \Hom_A( \tau_d^-\widetilde{ U },V^{ d } )
\]  
is surjective for each $\widetilde{ U } \in \cU$.  

\end{itemize}
To show this, observe that when $\delta$ is as in (i'), the contravariant defect functor $\delta^*$ is defined by the exact sequence
\[
  \Hom_A( -,V^{ d-1 } )
  \xrightarrow{}
  \Hom_A( -,V^d )
  \xrightarrow{}
  \delta^*( - )
  \xrightarrow{}
  0
\]
by \cite[def.\ 3.1]{Jasso-Kvamme_Introduction-to_higher_AR-theory}.  Hence the surjectivity required in (i') is equivalent to $\delta^*( \tau_d^-\cU ) = 0$, so (i') is equivalent to Lemma \ref{lem:splitting_vs_defects}, condition (iii).  This is again equivalent to Lemma \ref{lem:splitting_vs_defects}, condition (i), which is equal to condition (i) in the theorem.

(ii)$\Rightarrow$(i'):  This is clear.

(i')$\Leftrightarrow$(iii):  When $\delta$ is as in (i') and $\widetilde{ U } \in \cU$ is given, there is the following commutative diagram where $\alpha$, $\beta$, $\gamma$ are induced by $V^{ d-1 } \xrightarrow{} V^d$, and the vertical isomorphisms are given by Remark \ref{rmk:tau_d}(iii).
\[
\begin{tikzcd}
	{\Hom_A( \tau_d^-\widetilde{ U },V^{ d-1 } )} && {\Hom_A( \tau_d^-\widetilde{ U },V^{ d } )} \\
	{\underline{ \Hom }_A( \tau_d^-\widetilde{ U },V^{ d-1 } )} && {\underline{ \Hom }_A( \tau_d^-\widetilde{ U },V^{ d } )} \\
	{\overline{ \Hom }_A( \widetilde{ U },\tau_dV^{ d-1 } )} && {\overline{ \Hom }_A( \widetilde{ U },\tau_dV^{ d } )}
	\arrow["\alpha", from=1-1, to=1-3]
	\arrow[two heads, from=1-1, to=2-1]
	\arrow[two heads, from=1-3, to=2-3]
	\arrow["\beta", from=2-1, to=2-3]
	\arrow["{\rotatebox{270}{$\cong$}}", from=2-1, to=3-1]
	\arrow["{\rotatebox{270}{$\cong$}}", from=2-3, to=3-3]
	\arrow["\gamma"', from=3-1, to=3-3]
\end{tikzcd}
\]
We must show that $\alpha$ and $\gamma$ are surjective simultaneously.  This is equivalent to showing that $\alpha$ and $\beta$ are surjective simultaneously, which is true by Lemma \ref{lem:surjective_simultaneously}.
\end{proof}




\section{The spine of the lattice of $d$-torsion classes}
\label{sec:spine}


Recall from Remark \ref{rmk:cycles} the notion of a cycle in $\cM$.

\begin{Lemma}
\label{lem:deleting_an_indecomposable_from_a_set}
Assume that $\cM$ has no cycles.  

Let $\cX \subseteq \cM$ be a finite set of pairwise non-isomorphic indecomposable modules.  Then there exists $X \in \cX$ such that $X \in ( \cX \setminus X )^{ \perp }$.
\end{Lemma}

\begin{proof}
If the conclusion of the lemma were false, then each $X \in \cX$ would be outside $( \cX \setminus X )^{ \perp }$.  That is, for each $X \in \cX$ there would exist $X' \in \cX \setminus X$ and a homomorphism $X' \xrightarrow{} X$ that was not zero.  Since $X'$ and $X$ are different, hence non-isomorphic, the homomorphism would also not be an isomorphism.  Starting with an arbitrary $X_0 \in \cX\!$, we would hence be able to construct a sequence
\[
  \cdots \xrightarrow{} X_2 \xrightarrow{} X_1 \xrightarrow{} X_0
\]
with each $X_i$ in $\cX$ and each homomorphism not zero and not an isomorphism.  Since $\cX$ is finite, we would have $X_t = X_0$ for some $t \geqslant 1$, but this would contradict that $\cM$ has no cycles.
\end{proof}


\begin{Lemma}
\label{lem:deleting_an_indecomposable_from_a_d-torsion_class}
Assume that $\cM$ has no cycles.  

Let $\cV \subseteq \cW$ be a pair of splitting $d$-torsion classes differing by finitely many indecomposable modules.  

Pick a finite set $\cY \subseteq \cW \setminus \cV$ of pairwise non-isomorphic indecomposable modules such that $\add( \cY,\cV ) = \cW$.  Let $y \geqslant 0$ be the cardinality of $\cY\!$.

Then there exists a chain of splitting $d$-torsion classes in $\cM\!$,
\begin{equation}
\label{equ:lem:deleting_an_indecomposable_from_a_d-torsion_class:10}	
  \cV = \cW_y \subset \cdots \subset \cW_1 \subset \cW_0 = \cW\!,
\end{equation}
such that each is obtained by dropping one indecomposable module in $\cY$ from the class immediately above it.  That is, the inclusions are sharp, and there is a numbering $Y_1, \ldots, Y_y$ of the elements of $\cY$ such that $\add( Y_{ i+1 },\cW_{ i+1 } ) = \cW_i$ for $i \in \{\, 0, \ldots, y-1 \,\}$.
\end{Lemma}

\begin{proof}
Suppose that $\cW_t \subset \cdots \subset \cW_1 \subset \cW_0 = \cW$ and $Y_t, \ldots, Y_1$ have been constructed for some $t \geqslant 0$.  Since each $\cW_i$ is obtained by dropping the indecomposable module $Y_i$, we have
\begin{equation}
\label{equ:lem:deleting_an_indecomposable_from_a_d-torsion_class:20}
  \cW_t = \add( \cY \setminus \{\, Y_t, \ldots Y_1 \,\},\cV ).
\end{equation}
Note that if $t = 0$, then $\{\, Y_t, \ldots Y_1 \,\}$ is empty.

We claim that it is enough to assume $\cV \subset \cW_t$ and construct $\cW_{ t+1 }$ and $Y_{ t+1 }$.  Indeed, we can then iterate the construction, which at some point will give $\cV = \cW_t$ since $\cY$ is finite.  At this point, we will have constructed the chain \eqref{equ:lem:deleting_an_indecomposable_from_a_d-torsion_class:10}, and Equation \eqref{equ:lem:deleting_an_indecomposable_from_a_d-torsion_class:20} will imply $\cY \setminus \{\, Y_t, \ldots Y_1 \,\} = \emptyset$ whence $t = y$, so \eqref{equ:lem:deleting_an_indecomposable_from_a_d-torsion_class:10} has length $y$ as claimed.

So we assume $\cV \subset \cW_t$ and construct $\cW_{ t+1 }$ and $Y_{ t+1 }$.  Observe that $\cV \subset \cW_t$ combined with Equation \eqref{equ:lem:deleting_an_indecomposable_from_a_d-torsion_class:20} implies that $Y_t, \ldots, Y_1$ are not yet all the indecomposable modules in $\cY\!$.  Using Lemma \ref{lem:deleting_an_indecomposable_from_a_set}, pick $Y_{ t+1 } \in \cY \setminus \{\, Y_t, \ldots, Y_1 \,\}$ such that $Y_{ t+1 } \in ( \cY \setminus \{\, Y_{ t+1 }, \ldots, Y_1 \,\} )^{ \perp }$.  Define
\begin{equation}
\label{equ:lem:deleting_an_indecomposable_from_a_d-torsion_class:30}
  \cW_{ t+1 } = \add( \cY \setminus \{\, Y_{ t+1 }, \ldots, Y_1 \,\},\cV ).
\end{equation}

Equations \eqref{equ:lem:deleting_an_indecomposable_from_a_d-torsion_class:20} and \eqref{equ:lem:deleting_an_indecomposable_from_a_d-torsion_class:30} imply $\cW_{ t+1 } \subset \cW_t$ and $\add( Y_{ t+1 },\cW_{ t+1 } ) = \cW_t$, so to complete the proof, it is enough to show that $\cW_{ t+1 }$ is a splitting $d$-torsion class.  By Lemma \ref{lem:criterion_for_splitting}(iii), it suffices to let an indecomposable module $M \in \cM$ be given and prove $M \in \cW_{ t+1 }$ or $M \in \cW_{ t+1 }^{ \perp }$.  Suppose $M \not\in \cW_{ t+1 }$.  Then by Equations \eqref{equ:lem:deleting_an_indecomposable_from_a_d-torsion_class:20} and \eqref{equ:lem:deleting_an_indecomposable_from_a_d-torsion_class:30} there are two possibilities: $M \cong Y_{ t+1 }$ or $M \not\in \cW_t$.  

If $M \cong Y_{ t+1 }$ then we know $M \in ( \cY \setminus \{\, Y_{ t+1 }, \ldots, Y_1 \,\} )^{ \perp }$.  We also know $M \not\in \cV$ whence $M \in \cV^{ \perp }$ by Lemma \ref{lem:criterion_for_splitting}(iii) since $\cV$ is a splitting $d$-torsion class.  It now follows from Equation \eqref{equ:lem:deleting_an_indecomposable_from_a_d-torsion_class:30} that we have $M \in \cW_{ t+1 }^{ \perp }$.  If $M \not\in \cW_t$ then we have $M \in \cW_t^{ \perp }$ by Lemma \ref{lem:criterion_for_splitting}(iii) since $\cW_t$ is a splitting $d$-torsion class.  It now follows from $\cW_{ t+1 } \subset \cW_t$ that we have $M \in \cW_{ t+1 }^{ \perp }$.
\end{proof}


Recall that the spine of a lattice was originally defined for trim lattices by Thomas in \cite[p.\ 257]{Thomas_An-analogue-of-distributivity-for-ungraded-lattices}.  It consists of the elements which belong to a chain of maximal length.  


\begin{Theorem}
\label{thm:spine}
Assume that $\cM$ has no cycles and only finitely many isomorphism classes of indecomposable modules.  

Then the elements of the spine of the lattice of $d$-torsion classes of $\cM$ are precisely the splitting $d$-torsion classes.
\end{Theorem}

\begin{proof}
We let $m$ denote the number of isomorphism classes of indecomposable modules in $\cM$.

On the one hand, let $\cU \subseteq \cM$ be a splitting $d$-torsion class.  Applying Lemma \ref{lem:deleting_an_indecomposable_from_a_d-torsion_class} to the pairs $0 \subseteq \cU$ and $\cU \subseteq \cM$ produces two chains of splitting $d$-torsion classes in $\cM\!$.  Merging them gives a chain of splitting $d$-torsion classes in $\cM\!$,
\[
  0 \subset \cdots \subset \cU \subset \cdots \subset \cM\!,
\]
such that each is obtained by dropping one indecomposable module from the class immediately above it.  This property forces the chain to be of length $m$, and to be of maximal length in $\dtors\,\cM\!$, the lattice of $d$-torsion classes in $\cM$.  Since $\cU$ is in the chain, it is an element of the spine of $\dtors\,\cM$.  

Observe that as a byproduct, we learned that a chain of maximal length in $\dtors\,\cM$ has length $m$.

On the other hand, let $\cU \subseteq \cM$ be a $d$-torsion class which is an element of the spine of $\dtors\,\cM\!$.  Then there exists a chain of $d$-torsion classes in $\cM\!$,
\[
  0 = \cW_m \subset \cdots \subset \cW_1 \subset \cW_0 = \cM\!,
\]
which has maximal length, such that $\cU$ is one of the $\cW_i$.  By the previous part of the proof, the length of the chain is indeed $m$, and this implies that each class in the chain is obtained by dropping one indecomposable module from the class immediately above it.  To prove that $\cU$ is a splitting $d$-torsion class, it is enough to prove that so is each $\cW_i$.  Assume not.  Then it is easy to see that there exists $t$ such that $\cW_t$ is splitting but $\cW_{ t+1 }$ is not.  By Theorem \ref{thm:generalisation_of_ASS_VI1.7}(ii) there is an indecomposable module $W \in \cW_{ t+1 }$ such that
\begin{equation}
\label{equ:thm:spine:10}
  \tau_d^-W \not\in \cW_{ t+1 }.
\end{equation}
However, we know $\cW_{ t+1 } \subset \cW_t$, so $W$ is in the splitting $d$-torsion class $\cW_t$ whence Theorem \ref{thm:generalisation_of_ASS_VI1.7}(ii) implies 
\begin{equation}
\label{equ:thm:spine:20}
  \tau_d^-W \in \cW_t.
\end{equation}
There is an indecomposable module $Y \in \cM$ such that 
\begin{equation}
\label{equ:thm:spine:25}
  \add( Y,\cW_{ t+1 } ) = \cW_t,
\end{equation}
and Equations \eqref{equ:thm:spine:10} and \eqref{equ:thm:spine:20} imply $Y \cong \tau_d^-W$.  Hence the $d$-Auslander--Reiten sequence starting at $W$ is
\begin{equation}
\label{equ:thm:spine:30}
  0
  \xrightarrow{} W
  \xrightarrow{} M^1
  \xrightarrow{} \cdots
  \xrightarrow{} M^d
  \xrightarrow{} Y
  \xrightarrow{} 0.
\end{equation}
The sequence is minimal, see \cite[def.\ 3.1]{Iyama_Higher-dimensional-AR-theory} and \cite[rmk.\ 4.5]{Fedele_d-Auslander-Reiten-sequences_in_subcategories}, so since we have $W \in \cW_{ t+1 } \subset \cW_t$ and $Y \cong \tau_d^-W \in \cW_t$, we get $M^1 \in \cW_t$ since $\cW_t$ is closed under $d$-extensions; see \cite[lem.\ 3.8(1) and prop.\ 3.11]{August-et-al_A-characterisation-of-higher-torsion-classes}.  Now note that $Y$ cannot be a direct summand of $M^1$.  Indeed, \cite[lem.\ 3.8]{Pasquali_Tensor-products-of-n-complete-algebras} implies that each differential of \eqref{equ:thm:spine:30} is non-vanishing on each direct summand of an object, so if $Y$ were a direct summand of $M^1$, then \eqref{equ:thm:spine:30} would provide a cycle from $Y$ to $Y$ in $\cM$.  We have learned that $M^1 \in \cW_t$ and that $Y$ is not a direct summand of $M^1$.  By Equation \eqref{equ:thm:spine:25} this implies $M^1 \in \cW_{ t+1 }$.  But $\cW_{ t+1 }$ is closed under $d$-quotients by \cite[lem.\ 3.8(2) and prop.\ 3.11]{August-et-al_A-characterisation-of-higher-torsion-classes}, so we get $M^2, \ldots, M^d, Y \in \cW_{ t+1 }$.  By $Y \cong \tau_d^-W$ this contradicts Equation \eqref{equ:thm:spine:10}.
\end{proof}


\begin{Example}
[Dynkin type $A$]
\label{exa:Type_A}
Let $n \geqslant 1$ be an integer.  The $d$-Auslander algebra $A_n^d$ of Dynkin type $A_n$ was introduced in \cite[sec.\ 6.2]{Iyama_Cluster-tilting-for-higher-Auslander-algebras} based on the general theory developed in \cite{Iyama_Auslander-correspondence}.  There is a unique $d$-cluster tilting subcategory $\cM_n^d \subseteq \mod\,A_n^d$, whose lattice of $d$-torsion classes was investigated in \cite[sec.\ 5.5]{August-et-al_A-characterisation-of-higher-torsion-classes}.

The conditions of Theorem \ref{thm:spine} are satisfied by the concrete description of $\cM_n^d$ provided in \cite[thm./constr.\ 3.4 and thm.\ 3.6(3)]{Oppermann-Thomas_Higher-dimensional-cluster-combinatorics-and-representation-theory}.

Hence Theorem \ref{thm:spine} implies that the elements of the spine of the lattice of $d$-torsion classes of $\cM_n^d$ are precisely the splitting $d$-torsion classes.
\end{Example}


Note that the same method can also be applied to the higher Nakayama algebras of type $A$ introduced in \cite[def.\ 2.17]{Jasso-Kuelshammer-Psaroudakis-Kvamme_Higher_Nakayama_algebras_I} by Jasso, K\"{u}lshammer, Kvamme, and Psaroudakis.




\medskip
\noindent
{\bf Statement on AI.}
All results in the paper were originally proved by the authors, and all parts of the paper were written by them.  The proof of Theorem \ref{thm:generalisation_of_ASS_VI1.7} was improved by ideas supplied by Microsoft Copilot.  Substantial help with the references was provided by ChatGPT Pro 5.6.


\medskip
\noindent
{\bf Acknowledgement.}
This work was supported by an International Mobility Grant from the Research Council of Norway (project number 357034) and a NOVA Grant from Aarhus University Research Foundation (grant AUFF-E-2024-9-43).








\begin{thebibliography}{29}

%\bibitem{AI}  T.\ Aihara and O.\ Iyama, {\it Silting mutation in triangulated categories,} J.\ London Math.\ Soc.\ (2) {\bf 85} (2012), 633--668.

%\bibitem{A}  S.\ Al-Nofayee, {\it Simple objects in the heart of a t-structure,} J. Pure Appl.\ Algebra {\bf 213} (2009), 54--59.

\bibitem{Asadollahi-Jorgensen-Schroll-Treffinger_On-higher-torsion-classes} J.\ Asadollahi, P.\ J\o rgensen, S.\ Schroll, and H.\ Treffinger, {\it On higher torsion classes,} Nagoya Math.\ J.\ {\bf 248} (2022), 823--848.

\bibitem{Asai-Iyama-Mousavand-Paquette_Brick-splitting-torsion-pairs-and-trim-lattices}  S.\ Asai, O.\ Iyama, K.\ Mousavand, and C.\ Paquette, Brick-splitting torsion pairs and trim lattices, preprint (2025).  {\tt arXiv:2506.13602v1}

\bibitem{Assem-Simson-Skowronski_Elements-I}  I.\ Assem, D.\ Simson, and A.\ Skowro\'nski, ``Elements of the representation theory of associative algebras, Vol.\ 1, Techniques of representation theory'', London Math.\ Soc.\ Stud.\ Texts, Vol.\ 65, Cambridge University Press, Cambridge, 2006.

\bibitem{August-et-al_A-characterisation-of-higher-torsion-classes}  J.\ August, J.\ Haugland, K.\ M.\ Jacobsen, S.\ Kvamme, Y.\ Palu, and H.\ Treffinger, {\it A characterisation of higher torsion classes,} Forum Math.\ Sigma {\bf 13} (2025) e33 1--36.

\bibitem{August-et-al_Higher-torsion-classes-tau-d-tilting-theory-and-silting-complexes}  J.\ August, J.\ Haugland, K.\ M.\ Jacobsen, S.\ Kvamme, Y.\ Palu, and H.\ Treffinger, {\it Higher torsion classes, $\tau_d$-tilting theory and silting complexes,} preprint (2026).  {\tt arXiv:2602.03659v1}

%\bibitem{AQM} M.\ Auslander, ``Representation dimension of Artin algebras'', Queen Mary College Ma\-the\-ma\-tics Notes, Queen Mary College, London, 1971.  Reprinted pp.\ 505--574 in: ``Selected works of Maurice Auslander'', Vol.\ 1 (edited by Reiten, Smal\o, and Solberg), American Mathematical Society, Providence, 1999.

%\bibitem{AusRepI}  M.\ Auslander, {\it Representation theory of Artin algebras I,} Comm.\ Algebra {\bf 1} (1974), 177--268.

%\bibitem{AusRepII}  M.\ Auslander, {\it Representation theory of Artin algebras II}, Comm.\ Algebra {\bf 1} (1974), 269--310.

\bibitem{Auslander-Reiten-Smalo_Representation-theory-of-Artin-algebras}  M.\ Auslander, I.\ Reiten, and S.\ Smal\o, ``Representation theory of Artin algebras'', Cambridge Stud.\ Adv.\ Math., Vol.\ 36, Cambridge University Press, Cambridge, 1997.

%\bibitem{Auslander-Smalo_Preprojective-modules-over-Artin-algebras}  M.\ Auslander and S.\ O.\ Smal\o, {\it Preprojective modules over Artin algebras,} J.\ Algebra {\bf 66} (1980), 61--122.

%\bibitem{Avramov-Buchweitz-Iyengar}  L.\ L.\ Avramov, R.-O.\ Buchweitz, and S.\ Iyengar, {\it Class and rank of differential modules,} Invent.\ Math.\ {\bf 169} (2007), 1--35.

%\bibitem{AFH}  L.\ L.\ Avramov, H.-B.\ Foxby, and S.\ Halperin, Differential Graded homological algebra, unpublished notes.

%\bibitem{BBD}  A.\ A.\ Beilinson, J.\ Bernstein, and P.\ Deligne, {\it Faisceaux pervers,} Ast\'{e}risque {\bf 100} (1982) (proceedings of the conference ``Analysis and topology on singular spaces'', Vol.\ 1, Luminy, 1981).
%
%\bibitem{BR}  A.\ Beligiannis and I.\ Reiten, ``Homological and homotopical aspects of torsion theories'', Mem.\ Amer.\ Math.\ Soc., Vol.\ 188, 2007.

%\bibitem{Boekstedt-Neeman}  M.\ B\"{o}kstedt and A.\ Neeman, {\it Homotopy limits in triangulated categories,} Compositio Math.\ {\bf 86} (1993), 209--234. 

%\bibitem{BMRRT}  A.\ B.\ Buan, R.\ J.\ Marsh, M.\ Reineke, I.\ Reiten, and G.\ Todorov, {\it Tilting theory and cluster combinatorics,} Adv.\ Math.\ {\bf 204} (2006), 572--618.

%\bibitem{B}  T.\ B\"{u}hler, {\it Exact categories,} Expo.\ Math.\ {\bf 28} (2010), 1--69.

%\bibitem{Cartan-Eilenberg-Book}  H.\ Cartan and S.\ Eilenberg, ``Homological algebra'', Princeton University Press, Princeton, 1956.
%
%\bibitem{Chen-Iyengar}  X.-W.\ Chen and S.\ B.\ Iyengar, {\it Support and injective resolutions of complexes over commutative rings}, Homology Homotopy Appl.\ {\bf 12} (2010), 39--44.

\bibitem{Christensen-Foxby-Holm_Derived-category-methods}  L.\ W.\ Christensen, H.-B.\ Foxby, and H.\ Holm, ``Derived category methods in commutative algebra'', Springer Monographs in Mathematics, Springer, Cham, 2024.

%\bibitem{Christensen-Iyengar-Marley}  L.\ W.\ Christensen, S.\ B.\ Iyengar, and T.\ Marley, {\it Rigidity of Ext and Tor with coefficients in residue fields of a commutative Noetherian ring,} Proc.\ Edinb.\ Math.\ Soc.\ (2) {\bf 62} (2019), 305--321.

%\bibitem{CS2}  R.\ Coelho Sim\~{o}es, {\it  Hom-configurations and noncrossing partitions,} J.\ Algebraic Combin.\ {\bf 35} (2012), 313--343.
%
%\bibitem{CS3}  R.\ Coelho Sim\~{o}es, {\it Hom-configurations in triangulated categories generated by spherical objects,} J.\ Pure Appl.\ Algebra {\bf 219} (2015), 3322--3336.
%
%\bibitem{CS}  R.\ Coelho Sim\~{o}es, {\it Mutations of simple-minded systems in Calabi--Yau categories generated by a spherical object,} Forum Math.\ {\bf 29} (2017), 1065--1081.
%
%\bibitem{CSP1}  R.\ Coelho Sim\~{o}es and D.\ Pauksztello, {\it Simple-minded systems and reduction for negative Calabi--Yau triangulated categories,} Trans.\ Amer.\ Math.\ Soc.\ {\bf 373} (2020), 2463--2498.
%
%\bibitem{CSP2}  R.\ Coelho Sim\~{o}es and D.\ Pauksztello, {\it Torsion pairs in a triangulated category generated by a spherical object,} J.\ Algebra {\bf 448} (2016), 1--47.
%
%\bibitem{CSPP}  R.\ Coelho Sim\~{o}es, D.\ Pauksztello, and D.\ Ploog, Functorially finite hearts, simple-minded systems in negative cluster categories, and noncrossing partitions, preprint (2020).  {\tt arXiv:2004.00604v2}

%\bibitem{DSS}  I.\ Dell'Ambrogio, G.\ Stevenson, and J.\ \v{S}\v{t}ov\'{\i}\v{c}ek, {\it Gorenstein homological algebra and universal coefficient theorems,} Math.\ Z.\ {\bf 287} (2017), 1109--1155.

%\bibitem{Dickson}  S.\ E.\ Dickson, {\it A torsion theory for abelian categories,} Trans.\ Amer.\ Math.\ Soc.\ {\bf 121} (1966), 223--235.

%\bibitem{D2}  A.\ Dugas, {\it Tilting mutation of weakly symmetric algebras and stable equivalence,} Algebr.\ Represent.\ Theory {\bf 17} (2014), 863--884.
%
%\bibitem{D}  A.\ Dugas, {\it Torsion pairs and simple-minded systems in triangulated categories,} Appl.\ Categ.\ Structures {\bf 23} (2015), 507--526.

%\bibitem{Dyer}  M.\ Dyer, Exact subcategories of triangulated categories, preprint.  {\tt https://www3.nd.edu/\~{ }dyer/papers/extri.pdf}

%\bibitem{EEGR}  E.\ E.\ Enochs, S.\ Estrada, and J.\ R.\ Garc\'{\i}a-Rozas, {\it Gorenstein categories and Tate cohomology on projective schemes,} Math.\ Nachr.\ {\bf 281} (2008), 525--540.

%\bibitem{Enochs-Jenda-GorInjProj}  E.\ E.\ Enochs and O.\ M.\ G.\ Jenda, {\it Gorenstein injective and projective modules,} Math.\ Z.\ {\bf 220} (1995), 611--633.

%\bibitem{EJ}  E.\ E.\ Enochs and O.\ M.\ G.\ Jenda, ``Relative homological algebra'', de Gruyter Expositions in Mathematics, Vol.\ 30, Walter de Gruyter, Berlin New York, 2000.

%\bibitem{Enochs-Jenda-Xu}  E.\ E.\ Enochs, O.\ M.\ G.\ Jenda, and J.\ Xu, {\it Orthogonality in the category of complexes,} Math.\ J.\ Okayama Univ.\ {\bf 38} (1996), 25--46.

%\bibitem{Faith-Algebra-I}  C.\ Faith, ``Algebra: Rings, modules and categories I'', Grundlehren Math.\ Wiss., Vol.\ 190, Springer, Berlin Heidelberg New York, 1973.

%\bibitem{F}  F.\ Fedele, {\it Auslander--Reiten $(d+2)$-angles in subcategories and a $(d+2)$-angulated generalisation of a theorem by Br\"{u}ning,} J.\ Pure Appl.\ Algebra {\bf 223} (2019), 3554--3580.

\bibitem{Fedele_d-Auslander-Reiten-sequences_in_subcategories}  F.\ Fedele, {\it $d$-Auslander--Reiten sequences in subcategories,} Proc.\ Edinb.\ Math.\ Soc.\ (2) {\bf 63} (2020), 342--373.

\bibitem{Fedele-Jorgensen-Shah_The-index-in-d-exact-categories}  F.\ Fedele, P.\ J\o rgensen, and A.\ Shah, {\it The index in $d$-exact categories,} J.\ Algebra {\bf 686} (2026), 814--835.

%\bibitem{HBF-preprints-19}  H.-B.\ Foxby, A homological theory of complexes of modules, preprint (1981).
%
%\bibitem{HBF-Bounded-complexes}  H.-B.\ Foxby, {\it 
%Bounded complexes of flat modules,} J.\ Pure Appl.\ Algebra {\bf 15} (1979), 149--172.
%
%\bibitem{Freyd-book}  P.\ Freyd, ``Abelian categories'', Harper's Series in Modern Mathematics, 
%Harper \& Row, New York, 1964.

%\bibitem{Garcia-Rozas-book}  J.\ R.\ Garc\'{\i}a Rozas, ``Covers and envelopes in the category of complexes of modules'', Chapman \& Hall/CRC Res.\ Notes Math., Vol. 407, Chapman \& Hall/CRC, Boca Raton, Florida, 1999.

%\bibitem{Gillespie-book}  J.\ Gillespie, ``Abelian model category theory'', Cambridge Stud.\ Adv.\ Math., Vol.\ 215, Cambridge University Press, Cambridge, 2025.

%\bibitem{Gillespie-degreewise}  J.\ Gillespie, {\it Cotorsion pairs and degreewise homological model structures,} Homology Homotopy Appl.\ {\bf 10} (2008), 283--304.

%\bibitem{Gillespie-BLMS}  J.\ Gillespie, {\it Hereditary abelian model categories,} Bull.\ London Math.\ Soc.\ {\bf 48} (2016), 895--922.

%\bibitem{Gillespie-Exact}  J.\ Gillespie, {\it Model structures on exact categories,} J.\ Pure Appl.\ Algebra {\bf 215} (2011), 2892--2902.

%\bibitem{Goebel-Trlifaj-book}  R.\ G{\"o}bel and J.\ Trlifaj, ``Approximations and endomorphism algebras of modules'', de Gruyter Expositions in Mathematics, Vol.\ 41, Walter de Gruyter, Berlin, 2006.

%\bibitem{Goodearl-Warfield-book}  K.\ R.\ Goodearl and R.\ B.\ Warfield, ``An introduction to noncommutative Noetherian rings'', London Math.\ Soc.\ Stud.\ Texts, Vol.\ 61, Cambridge University Press, Cambridge, 2004.

%\bibitem{Grothendieck-Tohoku}  A.\ Grothendieck, {\it Sur quelques points d'alg\`{e}bre homologique,} Tohoku Math.\ J.\ (2) {\bf 9} (1957) 119--221.

%\bibitem{HRS}  D.\ Happel, I.\ Reiten, and S.\ Smal\o, ``Tilting in abelian categories and quasitilted algebras'', Mem.\ Amer.\ Math.\ Soc.\ {\bf 120} (1996), no.\ 575.

%\bibitem{Happel-book}  D.\ Happel, ``Triangulated categories in the representation theory of finite dimensional algebras'', London Math.\ Soc.\ Lecture Note Ser., Vol.\ 119, Cambridge University Press, Cambridge, 1988.
%
%\bibitem{Hilton-book}  P.\ Hilton, ``Homotopy theory and duality'', Gordon and Breach, New York--London--Paris, 1965.

%\bibitem{Hilton-Stammbach_Homological-algebra}  P.\ J.\ Hilton and U.\ Stammbach, ``A course in homological algebra'', Grad.\ Texts in Math., Vol.\ 4, Springer-Verlag, New York--Berlin, 1971.

%\bibitem{Holm-GHD}  H.\ Holm, {\it Gorenstein homological dimensions,} J.\ Pure Appl.\ Algebra {\bf 189} (2004), 167--193.
%
%\bibitem{HJ-Kyoto}  H.\ Holm and P.\ J\o rgensen, {\it Cotorsion pairs in categories of quiver representations,} Kyoto J.\ Math.\ {\bf 59} (2019), 575--606.
%
%\bibitem{HJ-Abel}  H.\ Holm and P.\ J\o rgensen, {\it A brief introduction to the $Q$-shaped derived category,} pp.\ 141--167 in ``Triangulated categories in representation theory and beyond'' (proceedings of the Abel Symposium 2022, edited by Bergh, Oppermann, and Solberg), Abel Symposia, Vol. 17, Springer, Cham, 2024.

%\bibitem{HJ-Model_cats}  H.\ Holm and P.\ J\o rgensen, {\it Model categories of quiver representations,} Adv.\ Math.\ {\bf 357} (2019), 106826. 

%\bibitem{HJ-JLMS}  H.\ Holm and P.\ J\o rgensen, {\it The $Q$-shaped derived category of a ring,} J.\ London Math.\ Soc.\ (2) {\bf 106} (2022), 3263--3316.
%
%\bibitem{HJ-TAMS}  H.\ Holm and P.\ J\o rgensen, {\it The $Q$-shaped derived category of a ring -- compact and perfect objects,} Trans.\ Amer.\ Math.\ Soc.\ {\bf 377} (2024), 3095--3128.

%\bibitem{HJY}  T.\ Holm, P.\ J\o rgensen, and D.\ Yang, {\it Sparseness of $t$-structures and negative Calabi--Yau dimension in triangulated categories generated by a spherical object,} Bull.\ London Math.\ Soc.\ {\bf 45} (2013), 120--130.

\bibitem{Hoshino_On-splitting-torsion-theories-induced-by-tilting-modules}  M.\ Hoshino, {\it On splitting torsion theories induced by tilting modules,} Comm.\ Algebra {\bf 11} (1983), 427--439.

%\bibitem{Hovey-CotMod}  M.\ Hovey, {\it Cotorsion pairs and model categories,} pp.\ 277--296 in ``Interactions between homotopy theory and algebra'' (Papers from the Summer School held at the University of Chicago, 26 July--6 August 2004. Edited by Avramov, Christensen, Dwyer, Mandell, and Shipley), Contemp.\ Math., Vol.\ 436, A\-me\-ri\-can Mathematical Society, Providence, RI, 2007.

%\bibitem{Hovey_MathZ}  M.\ Hovey, {\it Cotorsion pairs, model category structures, and representation theory,} Math.\ Z.\ {\bf 241} (2002), 553--592.

%\bibitem{Hovey_Book}  M.\ Hovey, ``Model categories'', Math.\ Surveys Monogr., Vol.\ 63, American Mathematical Society, Providence, 1999.

%\bibitem{Iacob-Iyengar}  A.\ Iacob and S.\ B.\ Iyengar, {\it Homological dimensions and regular rings,} J.\ Algebra {\bf 322} (2009), 3451--3458.

\bibitem{Iyama_Auslander-correspondence}  O.\ Iyama, {\it Auslander correspondence,} Adv.\ Math.\ {\bf 210} (2007), 51--82.

\bibitem{Iyama_Cluster-tilting-for-higher-Auslander-algebras}  O.\ Iyama, {\it Cluster tilting for higher Auslander algebras,} Adv.\ Math.\ {\bf 226} (2011), 1--61.

\bibitem{Iyama_Higher-dimensional-AR-theory}  O.\ Iyama, {\it Higher-dimensional Auslander-Reiten theory on maximal orthogonal subcategories,} Adv.\ Math.\ {\bf 210} (2007), 22--50.

%\bibitem{IJ}  O.\ Iyama and H.\ Jin, Positive Fuss--Catalan numbers and simple-minded systems in negative Calabi--Yau categories, preprint (2020).  {\tt arXiv:2002.09952v2}

%\bibitem{Iyama-Kato-Miyachi-N}  O.\ Iyama, K.\ Kato, and J. Miyachi, {\it Derived categories of $N$-complexes,} J.\ London Math.\ Soc.\ (2) {\bf 96} (2017), 687--716.
%
%\bibitem{Iyama-Minamoto-1}  O.\ Iyama and H.\ Minamoto, $\cA$-derived categories, in preparation.
%
%\bibitem{Iyama-Minamoto-2}  O.\ Iyama and H.\ Minamoto, {\it On a generalization of complexes and their derived categories,} extended abstract of a talk at the 47th Ring and Representation Theory Symposium, Osaka City University, 2014.

%\bibitem{IY}  O.\ Iyama and Y.\ Yoshino, {\it Mutation in triangulated categories and rigid Cohen-Macaulay mo\-du\-les,} Invent.\ Math.\ {\bf 172} (2008), 117--168.

%\bibitem{Iyengar-Krause-acyclicity}  S.\ Iyengar and H.\ Krause, {\it Acyclicity versus total acyclicity for complexes over noetherian rings,} Doc.\ Math.\ {\bf 11} (2006), 207--240.

\bibitem{Jasso_n-abelian}  G.\ Jasso, {\it $n$-abelian and $n$-exact categories}, Math.\ Z.\ {\bf 283} (2016), 703--759.

%\bibitem{Jasso-Q-shaped}  G.\ Jasso, {\it $Q$-shaped derived categories as derived categories of differential graded bimodules,} to appear in Annals of Representation Theory.

\bibitem{Jasso-Kuelshammer-Psaroudakis-Kvamme_Higher_Nakayama_algebras_I}  G.\ Jasso, J.\ K\"{u}lshammer, C.\ Psaroudakis, and S.\ Kvamme, {\it Higher Nakayama algebras I: Construction,} Adv.\ Math.\ {\bf 351} (2019), 1139--1200.

\bibitem{Jasso-Kvamme_Introduction-to_higher_AR-theory}  G.\ Jasso and S.\ Kvamme, {\it An introduction to higher Auslander--Reiten theory,} Bull.\ Lond.\ Math.\ Soc.\ {\bf 51} (2019), 1--24.

%\bibitem{Jin}  H.\ Jin, Reductions of triangulated categories and simple-minded collections, preprint (2019).  {\tt arXiv:1907.05114v2}
%
%\bibitem{JorARSub}  P.\ J\o rgensen, {\it Auslander--Reiten triangles in subcategories,} J.\ K-theory {\bf 3} (2009), 583--601.

%\bibitem{JorSMS}  P.\ J\o rgensen, Abelian subcategories of triangulated categories induced by simple minded systems, preprint (2020).  {\tt arXiv:2010.11799v2}

\bibitem{Jorgensen_Torsion-classes-in-higher-homological-algebra}  P.\ J\o rgensen, {\it  Torsion classes and t-structures in higher homological algebra,} Internat.\ Math.\ Res.\ Notices {\bf 2016} (2016), 3880--3905.

%\bibitem{Kapranov}  M.\ M.\ Kapranov, On the $q$-analog of homological algebra, preprint (1996).  {\tt arXiv:9611005v1.}

%\bibitem{Kashiwara-Schapira}  M.\ Kashiwara and P.\ Schapira, ``Categories and sheaves'', Grundlehren Math.\ Wiss., Vol.\ 332, Springer, Berlin Heidelberg, 2006.

%\bibitem{K}  B.\ Keller, {\it Chain complexes and stable categories,} Manuscripta Math.\ {\bf 67} (1990), 379--417.
%
%\bibitem{KellerOrbit}  B.\ Keller, {\it On triangulated orbit categories,} Doc.\ Math.\ {\bf 10} (2005), 551--581.

%\bibitem{Klapproth}  C.\ Klapproth, $n$-exact categories arising from $( n+2 )$-angulated categories, preprint (2021).  {\tt arXiv:2108.04596v2}

%\bibitem{KL}  S.\ K\"{o}nig and Y.\ Liu, {\it Simple-minded systems in stable module categories,} Quart.\ J.\ Math.\ Oxford {\bf 63} (2012), 653--674.
%
%\bibitem{KY}  S.\ K\"{o}nig and D.\ Yang, {\it Silting objects, simple-minded collections, t-structures and co-t-structures for finite-dimensional algebras,} Doc.\ Math.\ {\bf 19} (2014), 403--438.
%
%\bibitem{Krause}  H.\ Krause, {\it Krull--Schmidt categories and projective covers,} Expo.\ Math.\ {\bf 33} (2015), 535--549.

%\bibitem{Krause}  H.\ Krause, {\it The stable derived category of a noetherian scheme,} Compos.\ Math.\ {\bf 141} (2005), 1128--1162.

\bibitem{Kvamme_Higher-extension-closure-and-d-exact-categories}  S.\ Kvamme, {\it Higher extension closure and $d$-exact categories,} Adv.\ Math.\ {\bf 499} (2026), 111054.

%\bibitem{L}  M.\ Linckelmann, On abelian subcategories of triangulated categories, preprint (2021).  {\tt arXiv:2008.03199v2}

%\bibitem{MacLane-book}  S.\ Mac Lane, ``Categories for the working mathematician'', second edition, Grad.\ Texts in Math., Vol.\ 5, Springer, New York, 1998.
%
%\bibitem{Mangeney-Peskine-Szpiro}  M.\ Mangeney, C.\ Peskine, and L.\ Szpiro, {\it Anneaux de Gorenstein, et torsion en alg\`{e}bre commutative,} notes based on lectures by M.\ Auslander, pp.\ 2--69 in: S\'{e}minaire Samuel. Alg\`{e}bre commutative, tome 1 (1966-1967).  

%\bibitem{Miyashita_Tilting-modules-of-finite-projective-dimension}  Y.\ Miyashita, {\it Tilting modules of finite projective dimension,} Math.\ Z.\ {\bf 193} (1966), 113--146.

%\bibitem{NP}  H.\ Nakaoka and Y.\ Palu, {\it Extriangulated categories, Hovey twin cotorsion pairs and model structures,} Cahiers Topologie G\'{e}om. Diff\'{e}rentielle Cat\'{e}g., {\bf 60} (2019), 117--193.

%\bibitem{N}  A.\ Neeman, ``Triangulated categories'', Ann.\ of Math.\ Stud., Vol.\ 148, Princeton University Press, Princeton,
%2001. 

%\bibitem{Nkansah-differential-modules}  D.\ Nkansah, Differential modules: A perspective on Bass' question, preprint (2025).  {\tt arXiv:2504.15981v1}
%
%\bibitem{Oberst-Roehrl}  U.\ Oberst and H.\ Rohrl, {\it Flat and coherent functors,} J.\ Algebra {\bf 14} (1970), 91--105.

\bibitem{Oppermann-Thomas_Higher-dimensional-cluster-combinatorics-and-representation-theory}  S.\ Oppermann and H.\ Thomas, {\it Higher-dimensional cluster combinatorics and representation theory}, J.\ Eur.\ Math.\ Soc.\ (JEMS) {\bf 14} (2012), 1679--1737.


\bibitem{Pasquali_Tensor-products-of-n-complete-algebras}  A.\ Pasquali, {\it Tensor products of $n$-complete algebras,} J.\ Pure Appl.\ Algebra {\bf 223} (2019), 3537--3553.

%\bibitem{Q}  D.\ Quillen, {\it Higher algebraic K-theory: I,} pp.\ 85--147 in: ``Algebraic K-theory, I: Higher K-theories'' (proceedings of the conference held at the Seattle Research Center of the Battelle Memorial Institute, 28 August to 8 September 1972, edited by Bass), Lecture Notes in Math., Vol.\ 341, Springer-Verlag, Berlin-New York, 1973.
%
%\bibitem{Rickard}  J.\ Rickard, {\it Equivalences of derived categories for symmetric algebras,} J.\ Algebra {\bf 257} (2002), 460--481.
%
%\bibitem{R}  C.\ M.\ Ringel, {\it Appendix: Some remarks concerning tilting modules and tilted algebras. Origin. Relevance. Future.,} pp.\ 413--472 in ``Handbook of tilting theory'' (edited by Angeleri-H\"{u}gel, Happel, and Krause), London Math.\ Soc.\ Lecture Note Ser., Vol.\ 332, Cambridge University Press, Cambridge, 2007.

%\bibitem{Ringel-Zhang} C.\ M.\ Ringel and P.\ Zhang, {\it Representations of quivers over the algebra of dual numbers,} J.\ Algebra {\bf 475} (2017), 327--360.
%
%\bibitem{Roig}  A.\ Roig, {\it Minimal resolutions and other minimal models,} Publ.\ Mat.\ {\bf 37} (1993), 285--303.

%\bibitem{Sch}  O.\ M.\ Schn\"{u}rer, Simple-minded subcategories and t-structures, preprint (2013).

%\bibitem{Schiffler}  R.\ Schiffler, ``Quiver representations'', CMS Books Math., Springer, Cham, 2014.

\bibitem{Segovia_P-comma-phi-Tamari-and-higher-torsion-lattices-of-type-A}  A.\ Segovia, $( P,\phi )$-Tamari and higher torsion lattices of type $A$, preprint (2026).  {\tt arXiv:2606.00273v1}

\bibitem{Segovia_P-comma-phi-Tamari-lattices}  A.\ Segovia, {\it $( P,\phi )$-Tamari lattices,} S\'{e}m.\ Lothar.\ Combin.\ {\bf 93B} (2025), article \#71.

%\bibitem{Slaftsos-Vitoria}  A.\ Slaftsos and J.\ Vit\'{o}ria, Relative $Q$-shaped homological algebra, preprint (2026).  {\tt arxiv:2602.22986v1.}
%
%\bibitem{Stenstroem}  B.\ Stenstr\"{o}m, ``Rings of quotients'', Grundlehren Math.\ Wiss., Vol.\ 217, Springer, Berlin--Heidelberg--New York, 1975.

\bibitem{Thomas_An-analogue-of-distributivity-for-ungraded-lattices}  H.\ Thomas, {\it An analogue of distributivity for ungraded lattices,} Order {\bf 23} (2006), 249--269.

%\bibitem{Wei}  J. Wei, {\it Gorenstein homological theory for differential modules,} Proc.\ Roy.\ Soc.\ Edinburgh Sect.\ A {\bf 145} (2015), 639--655.

%\bibitem{W}  J.\ Woolf, {\it Stability conditions, torsion theories and tilting,} J.\ London Math.\ Soc.\ (2) {\bf 82} (2010), 663--682.

%\bibitem{Xu-Flat-covers}  J.\ Xu, ``Flat covers of modules'', Lecture Notes in Math., Vol.\ 1634, Springer, New York, 1996.

\end{thebibliography}

\end{document}













